% !TeX root = ../../Book4.tex
% 纲要标题：## 第1章　复分析的几何方法
% 纲要位置：几何学计划纲要.md，第 3477 行。
% 正文状态：初稿；尚未正式审核。

% 纲要内容开始
% 
% > 本章目标与证明范围：面向研究生重建几何所需的一元复分析，重心是局部映射、解析延拓、紧性和极值构造，而不是把本科章节依次缩写。
% 
% 纲要内容结束

\chapter{复分析的几何方法}
\label{g4:ch:01}
\BookChapterNumber{g4:ch:01}

\begin{chapterabstract}
	一个复导数的存在会同时限制函数的局部形状、零点分布和整体延拓。本章从复可微与积分公式建立解析性，证明非恒定全纯函数的局部幂模型；随后以留数、绕数和辐角原理计数原像，用函数芽与单值定理描述沿路径的延拓。极值原理及圆盘上的收缩性质给出刚性和归一化唯一性，正规族则把有界函数列转化为可用的全纯极限。最后通过导数极值构造，证明每个单连通真平面域都与单位圆盘双全纯等价，并说明域内等价与边界对应各自包含什么信息。
\end{chapterabstract}

\begin{learninggoals}
	\begin{itemize}
		\item 从复差商得到 Cauchy--Riemann 方程，用三角形积分和 Cauchy 公式证明全纯函数的幂级数表示。
		\item 计算零点阶数，把非恒定全纯函数在一点附近化为幂映射，并判断逆分支的存在范围。
		\item 用主部、留数、辐角原理和 Rouché 定理计算零极点与原像重数，辨明边界和内部假设的作用。
		\item 用有限延拓链追踪函数芽，证明同伦不变性，并由单值定理构造全纯对数和根。
		\item 推导 Schwarz--Pick 的两点及微分不等式，解释等号如何确定圆盘自同构。
		\item 以局部有界性取得全纯极限，完成 Riemann 映照定理的极值证明，并判断哪些边界结论需要额外条件。
	\end{itemize}
\end{learninggoals}

方程 \(w^2=z\) 在每个非零点附近都能解出两个全纯分支，却不能在整个穿孔平面上选择一个连续的平方根。另一方面，一块边界形状很不规则的单连通平面区域，仍可能拥有一个把它整体变成圆盘的复坐标。前一个问题要求追踪局部逆沿环路怎样返回，后一个问题要求从一族候选坐标中找出真正覆盖圆盘的映射。一元复分析提供了处理这两种问题的共同工具：积分把局部微分条件传到边界，整数重数在扰动中保持，紧性使一列越来越好的候选映射产生极限。

\ProseParagraphBreak

这些问题也解释了为何复曲线的理论要先建立平面函数的几何。将来在曲面上更换全纯坐标时，导数非零保证坐标之间可逆；允许分歧时，局部坐标中的幂指数给出重数；在不同图中继续一个函数时，还要核对它绕行之后是否回到同一个函数芽。本章在平面区域中逐项完成这些论证，为下一章的复图册、覆盖与分歧准备可直接使用的结果。

\ProseParagraphBreak

本章需要复数的代数、实多元微积分、幂级数的一致收敛，以及开集、紧集和连通性的基本性质。全纯逆函数、同伦下的解析延拓和正规族紧性均在本章证明。带孔区域的积分使用平面 Green 公式，调用时恢复其假设、方向及第一卷中的证明位置；读者掌握该积分公式即可从这里开始复几何路线。第三卷的函数空间与椭圆理论在本章尚无必要。复积分与局部映射可对照\cite[第4章第2--3节]{Ahlfors1979Complex}，正规族可查\cite[第5章第5.3--5.4节]{Ahlfors1979Complex}；映照构造与解析延拓分别见\cite[第6章第1.1节; 第8章第1.4--1.6节]{Ahlfors1979Complex}。

% !TeX root = ../../Book4.tex
% 纲要标题：### 1.1 全纯性和局部映射
% 纲要位置：几何学计划纲要.md，第 3481 行。
% 正文状态：初稿；尚未正式审核。

% 纲要内容开始
% 
% * Cauchy–Riemann 方程与幂级数
% * Cauchy 积分及解析性的证明
% * 零点阶数与局部幂模型
% * 全纯逆函数及分支
% 
% 纲要内容结束

\section{全纯性和局部映射}
\label{g4:sec:0101}

映射 \(z\mapsto z^2\) 把辐角加倍，却在每个非零点附近保持相交曲线之间的有向角。这两个事实分别涉及有限范围内的绕转和一点处的线性近似。全纯函数的几何从这种区别开始：复导数约束无穷小形状，幂级数则把临界点附近的整个映射化为一个幂映射。我们先从复可微出发证明幂级数的存在，再建立这一局部模型。

\ProseParagraphBreak

本章用 \(D(a,r)=\{z\in\mathbb C:|z-a|<r\}\) 表示开圆盘，\(\mathbb D=D(0,1)\) 表示单位圆盘。\term[quyu]{区域}[domain]指非空连通开集；开集未必连通。圆周的正向为逆时针，复平面的定向由有序实坐标 \((x,y)\) 给定。

\subsection{Cauchy–Riemann 方程与幂级数}
\label{g4:subsec:0101:01}
把 \(f=u+\mathrm i v\) 看成实平面的映射时，它的一阶变化通常由一个任意实矩阵给出。复可微要求这个矩阵等于乘以某个复数；非零复数的乘法是一次旋转与正比例伸缩的复合，因而把两个方向之间的角度保留下来。

\begin{definition}
	\label{g4:def:0101:holomorphic}
	设 \(U\subset\mathbb C\) 开，\(a\in U\)。若极限
	\[
		f'(a)=\lim_{h\to0,\ h\ne0}\frac{f(a+h)-f(a)}{h}
	\]
	存在，则称 \(f:U\to\mathbb C\) 在 \(a\) 处复可微；这里 \(h\) 沿复平面中的任意方向趋于零。若 \(f\) 在 \(U\) 每点复可微，则称它在 \(U\) 上\term[quanchun]{全纯}[holomorphic]，全体这样的函数记为 \(\mathcal O(U)\)。在整个 \(\mathbb C\) 上全纯的函数称为\term[zhenghanshu]{整函数}[entire function]。
	\symidx{\(\mathcal O(U)\)：开集 \(U\) 上的全纯函数}
\end{definition}

复导数存在等价于 \(f(a+h)=f(a)+Ah+o(|h|)\)，其中 \(A=f'(a)\)。因此复可微已经蕴含实可微。反过来，若 \(f=u+\mathrm i v\) 在 \(a\) 实可微，则
\[
	\mathrm{D}f_a=
	\begin{pmatrix}u_x&u_y\\v_x&v_y\end{pmatrix}_a
\]
表示复数乘法，当且仅当 \(u_x=v_y\)、\(u_y=-v_x\)。这就是\term[cauchyriemannfangcheng]{Cauchy--Riemann 方程}[Cauchy--Riemann equations]。\footnote{Cauchy，Augustin-Louis Cauchy（1789--1857），法国数学家，研究实分析、复分析与置换群。Riemann，Georg Friedrich Bernhard Riemann（1826--1866），德国数学家，建立复函数的曲面观点，并深刻影响了微分几何与积分理论。}在此情形
\begin{equation}
	\label{g4:eq:0101:complex-linear}
	f'=u_x+\mathrm i v_x,\qquad
	\mathrm{D}f=
	\begin{pmatrix}\operatorname{Re}f'&-\operatorname{Im}f'\\
	\operatorname{Im}f'&\operatorname{Re}f'\end{pmatrix},\qquad
	\det\mathrm{D}f=|f'|^2.
\end{equation}
尤其对 \(C^1\) 函数，满足这些方程与全纯等价。后面还要证明，全纯函数自动具有任意阶连续导数，因而定义中无需另加 \(C^1\) 假设。

\ProseParagraphBreak

只检查一点处的偏导关系并不足够。例如令 \(f(0)=0\)，在 \((x,y)\ne(0,0)\) 时令
\[
	f(x+\mathrm i y)=\frac{x^3y}{x^4+y^2}.
\]
它在原点连续，因为 \(|x^3y|/(x^4+y^2)\le |x|/2\)；沿两条坐标轴为零，所以四个实偏导在原点满足上述方程。但沿 \(y=x^2\)，有 \(f(x+\mathrm i x^2)=x/2\)，故复差商趋于 \(1/2\)，而沿实轴差商为零。这里缺少的正是实可微性。

\begin{definition}
	\label{g4:def:0101:analytic}
	若函数在每点 \(a\) 的某圆盘内都可表示为收敛幂级数 \(\sum_{n\ge0}c_n(z-a)^n\)，则称它\term[jiexihanshu]{解析}[analytic]。
\end{definition}

收敛半径为 \(R>0\) 的幂级数在 \(|z-a|\le r<R\) 上一致绝对收敛，且可逐项求导。具体说，选 \(r<\rho<R\)，则 \(|c_n|\rho^n\) 有界，而 \(\sum n(r/\rho)^{n-1}\) 收敛，控制了导数级数的一致收敛；积分基本定理随后说明其和确为原函数的导数。逐次使用这一估计，得到所有阶导数。因此解析蕴含全纯，且 \(c_n=f^{(n)}(a)/n!\)。困难的反向结论是：仅有一次复导数，为何就能恢复整个幂级数？

\ProseParagraphBreak

全纯函数的和、积、复合仍全纯；分母非零处的商也全纯。这些性质由差商法则直接得到。特别地，\(\exp z=\sum z^n/n!\) 全纯，满足 \((\exp z)'=\exp z\) 和 \(\exp(z+w)=\exp z\exp w\)；后者可由绝对收敛级数的乘积逐项合并得到。取 \(w=-z\) 可知指数从不为零。它将成为构造局部对数和根的工具。

\subsection{Cauchy 积分及解析性的证明}
\label{g4:subsec:0101:02}
边界上的函数值是否足以决定内部？实可微函数可以在圆盘内部任意作紧支扰动；全纯函数的答案却完全不同。关键是先证明小三角形边界上的积分为零，再把这个消去关系用于带一个极点的积分核。

\ProseParagraphBreak

对分片 \(C^1\) 曲线 \(\gamma:[\alpha,\beta]\to U\) 及连续函数 \(f\)，定义
\[
	\int_\gamma f(z)\,\mathrm{d}z
	=\int_\alpha^\beta f(\gamma(t))\gamma'(t)\,\mathrm{d}t.
\]
积分按分段相加，反向改变符号，且绝对值不超过曲线长度乘以曲线上 \(|f|\) 的上界。若 \(F'=f\)，链式法则给出 \(\int_\gamma f\,\mathrm{d}z=F(\gamma(\beta))-F(\gamma(\alpha))\)，所以原函数的存在会使闭曲线积分消失。

\begin{theorem}[Cauchy--Goursat 定理]
	\label{g4:thm:0101:goursat}
	设 \(U\subset\mathbb C\) 开，\(f\in\mathcal O(U)\)。若闭三角形 \(T\) 包含于 \(U\)，则
	\[
		\int_{\partial T}f(z)\,\mathrm{d}z=0.
	\]
\end{theorem}
这称为\term[cauchygoursatdingli]{Cauchy--Goursat 定理}[Cauchy--Goursat theorem]。\footnote{Goursat，Édouard Jean-Baptiste Goursat（1858--1936），法国数学家，研究复分析与微分方程；其积分定理去掉了对全纯函数导数连续性的预设。}
\begin{proof}
	记原积分为 \(I\)。连接三条边的中点，得到四个相似小三角形；内部边两两抵消，故至少一个小三角形的积分绝对值不小于 \(|I|/4\)。重复选取，得到套叠闭三角形 \(T_n\)，直径和周长分别为 \(2^{-n}\operatorname{diam}T\)、\(2^{-n}\ell(\partial T)\)，且其积分绝对值不小于 \(4^{-n}|I|\)。紧性及直径趋零保证交集为一点 \(p\in U\)。

	在 \(p\) 复可微给出
	\[
		f(z)=f(p)+f'(p)(z-p)+(z-p)\varepsilon(z),\qquad \varepsilon(z)\longrightarrow0.
	\]
	常数项和一次项有显式原函数，闭边界积分为零。因此
	\[
		4^{-n}|I|\le 4^{-n}\ell(\partial T)\operatorname{diam}T
		\sup_{z\in T_n}|\varepsilon(z)|.
	\]
	约去 \(4^{-n}\) 并令 \(n\to\infty\)，得到 \(I=0\)。
\end{proof}

\begin{lemma}
	\label{g4:lem:0101:primitive}
	设 \(D\subset\mathbb C\) 是开圆盘，\(g:D\to\mathbb C\) 连续。若 \(g\) 在每个闭三角形 \(T\subset D\) 的边界上积分为零，则 \(g\) 在 \(D\) 上有全纯原函数。特别地，全纯函数在圆盘上有原函数。
\end{lemma}
\begin{proof}
	固定 \(b\in D\)，沿直线段定义 \(G(z)=\int_{[b,z]}g(\zeta)\,\mathrm{d}\zeta\)。圆盘凸，三角形的零积分给出
	\[
		G(z+h)-G(z)=\int_{[z,z+h]}g(\zeta)\,\mathrm{d}\zeta
		=h\int_0^1g(z+th)\,\mathrm{d}t.
	\]
	连续性使右端除以 \(h\) 后趋于 \(g(z)\)。于是 \(G'=g\)，原函数在整个圆盘上定义。
\end{proof}

下一步会用到一个轻微推广：若 \(g\) 在圆盘上连续，仅在一点 \(p\) 可能不全纯，上述三角形积分仍为零。若三角形遇到 \(p\)，先把它分成以 \(p\) 为顶点的至多三个三角形，再在该顶点切去直径 \(\varepsilon\) 的小三角形。剩余多边形可剖分成不经过 \(p\) 的三角形，对它们用\cref{g4:thm:0101:goursat}。原积分与剩余积分的差由总长度 \(O(\varepsilon)\) 的线段积分组成；\(g\) 在附近有界，故差趋零。点落在边上时用同样切割。这样，\cref{g4:lem:0101:primitive}也适用于这个连续延拓。

\begin{theorem}[Cauchy 积分公式]
	\label{g4:thm:0101:cauchy-formula}
	设 \(U\subset\mathbb C\) 开，\(f\in\mathcal O(U)\)，且 \(\overline{D(a,r)}\subset U\)。对每个 \(z\in D(a,r)\)，有
	\begin{equation}
		\label{g4:eq:0101:cauchy-formula}
		f(z)=\frac{1}{2\pi\mathrm i}\int_{|\zeta-a|=r}
		\frac{f(\zeta)}{\zeta-z}\,\mathrm{d}\zeta.
	\end{equation}
	进而 \(f\) 在该圆盘上有收敛幂级数
	\begin{equation}
		\label{g4:eq:0101:power-series}
		f(z)=\sum_{n=0}^\infty c_n(z-a)^n,\qquad
		c_n=\frac{1}{2\pi\mathrm i}\int_{|\zeta-a|=r}
		\frac{f(\zeta)}{(\zeta-a)^{n+1}}\,\mathrm{d}\zeta
		=\frac{f^{(n)}(a)}{n!}.
	\end{equation}
	级数在每个较小闭圆盘上一致绝对收敛。因此全纯与解析等价，全纯函数为 \(C^\infty\)。
\end{theorem}
公式\eqref{g4:eq:0101:cauchy-formula}称为\term[cauchyjifengongshi]{Cauchy 积分公式}[Cauchy integral formula]。
\begin{proof}
	因闭圆盘紧含于 \(U\)，可把半径略增，仍留在 \(U\)。对固定的内部点 \(z\)，在稍大的圆盘上令
	\[
		g(\zeta)=\frac{f(\zeta)-f(z)}{\zeta-z}\quad(\zeta\ne z),\qquad g(z)=f'(z).
	\]
	它在 \(z\) 连续，在其余点全纯。前面的推广及原函数引理说明其圆周积分为零。另一方面，因 \(|z-a|<r\)，在积分圆周上可以一致展开
	\[
		\frac1{\zeta-z}
		=\sum_{n=0}^\infty\frac{(z-a)^n}{(\zeta-a)^{n+1}}.
	\]
	参数化 \(\zeta=a+r e^{\mathrm i t}\) 可见，右端各项积分除 \(n=0\) 给出 \(2\pi\mathrm i\) 外，其余均为零。因此 \(\int(\zeta-z)^{-1}\,\mathrm{d}\zeta=2\pi\mathrm i\)。把 \(g\) 的零积分展开，便得积分公式。

	将同一几何级数代入积分公式。在 \(|z-a|\le s<r\) 上，\(|z-a|/r\le s/r<1\)，所以可一致逐项积分，得到所列 \(c_n\)。若圆周上的最大模为 \(M_r\)，则 \(|c_n|\le M_r/r^n\)，再次保证较小闭圆盘上的一致绝对收敛。幂级数逐项求导已在前面证明，故 \(c_n=f^{(n)}(a)/n!\)，并得到所有阶导数。
\end{proof}

\begin{proposition}[Cauchy 估计]
	\label{g4:prop:0101:cauchy-estimate}
	设 \(f\) 在闭圆盘 \(\overline{D(a,r)}\) 的邻域全纯，\(M_r=\max_{|z-a|=r}|f(z)|\)。则对所有整数 \(n\ge0\)，
	\[
		|f^{(n)}(a)|\le\frac{n!M_r}{r^n}.
	\]
	更一般地，设 \(U\subset\mathbb C\) 开，\(f\in\mathcal O(U)\)，\(K\subset U\) 紧。若 \(\delta>0\)，所有 \(D(z,\delta)\)（\(z\in K\)）的闭包都包含于 \(U\)，且 \(|f|\le M\) 在这些圆盘上成立，则 \(\sup_K|f^{(n)}|\le n!M/\delta^n\)。
\end{proposition}
\begin{proof}
	对\eqref{g4:eq:0101:power-series}中的系数积分取绝对值，使用圆周长度 \(2\pi r\) 即得第一式。对每个 \(z\in K\) 以它为圆心重复同一估计，便得第二式。
\end{proof}

这个估计把函数值的控制转成所有阶导数的控制。例如有界整函数 \(f\) 若满足 \(|f|\le M\)，则对任意 \(a\) 和任意 \(r>0\)，有 \(|f'(a)|\le M/r\)。令 \(r\to\infty\)，得到 \(f'=0\)，再沿线段积分可知 \(f\) 为常数。这里半径能够无限增大，正是定义域为整个平面的作用。

\ProseParagraphBreak

有了光滑性，就能把圆盘公式扩展到带孔区域。所需实变量公式为
\begin{equation}
	\label{g4:eq:0101:green}
	\int_{\partial G}(P\,\mathrm{d}x+Q\,\mathrm{d}y)
	=\int_G(Q_x-P_y)\,\mathrm{d}x\,\mathrm{d}y.
\end{equation}
这里 \(G\) 有界，边界由有限条互不相交的分片 \(C^2\) 正则简单闭曲线组成，角点内角在 \((0,2\pi)\) 内，\(P,Q\) 在 \(\overline G\) 的邻域为 \(C^1\)；边界按区域位于左侧定向。第一卷第9.1节引理9.1.9 已由积分基本定理、有限局部分解和角点截去极限证明这一 Green 公式。\footnote{Green，George Green（1793--1841），英国数学家，以位势理论、Green 函数和积分公式著称。}复系数情形分别用于实部和虚部即可。本章对一般边界区域均使用上述条件；圆盘和圆环是最常用的情形。

\begin{proposition}
	\label{g4:prop:0101:cauchy-domain}
	设 \(G\subset\mathbb C\) 是有界区域，边界满足公式\eqref{g4:eq:0101:green}所述条件，\(f\) 在 \(\overline G\) 的邻域全纯。则
	\[
		\int_{\partial G} f(\zeta)\,\mathrm{d}\zeta=0,
		\qquad
		f(z)=\frac1{2\pi\mathrm i}\int_{\partial G}
		\frac{f(\zeta)}{\zeta-z}\,\mathrm{d}\zeta\quad(z\in G).
	\]
	有孔时，外边界逆时针，内边界顺时针。
\end{proposition}
\begin{proof}
	写 \(f\,\mathrm{d}z=(u\,\mathrm{d}x-v\,\mathrm{d}y)+\mathrm i(v\,\mathrm{d}x+u\,\mathrm{d}y)\)。由\eqref{g4:eq:0101:green}，两部分的面积被积函数分别为 \(-v_x-u_y\) 与 \(u_x-v_y\)，恰由 Cauchy--Riemann 方程为零。使用该公式合法，因为\cref{g4:thm:0101:cauchy-formula}已经给出 \(f\) 的光滑性。

	对第二式，在 \(z\) 周围删去闭小圆盘 \(\overline{D(z,\varepsilon)}\subset G\)。核 \(f(\zeta)/(\zeta-z)\) 在剩余区域的闭包附近全纯，第一式说明其原边界积分等于小圆的逆时针积分。后者写成
	\[
		\mathrm i\int_0^{2\pi} f(z+\varepsilon e^{\mathrm i t})\,\mathrm{d}t
	\]
	并令 \(\varepsilon\to0\)，即得到 \(2\pi\mathrm i f(z)\)。
\end{proof}

\subsection{零点阶数与局部幂模型}
\label{g4:subsec:0101:03}
幂级数的第一个非零项不仅给出近似，还决定零点的重数。把其后的非零因子吸收到新坐标里，便能把近似变成精确等式。

\begin{proposition}
	\label{g4:prop:0101:zero-order}
	设 \(f\) 在点 \(a\) 的邻域全纯且 \(f(a)=0\)。恰有下列一种情形成立：\(f\) 在 \(a\) 的某邻域恒为零；或者存在唯一整数 \(m\ge1\) 及全纯函数 \(u\)，使
	\[
		f(z)=(z-a)^m u(z),\qquad u(a)\ne0.
	\]
	后一情形的 \(a\) 是孤立零点，整数 \(m\) 称为其\term[lingdiandejieshu]{零点阶数}[order of a zero]或重数，记为 \(\operatorname{ord}_a f\)。
\end{proposition}
\begin{proof}
	在 \(a\) 展开幂级数。若全部系数为零，则函数在收敛圆盘上为零；否则取第一个非零系数的指标 \(m\)，提出 \((z-a)^m\) 即得 \(u\)。其常数项非零，连续性保证缩小圆盘后 \(u\) 处处非零。幂级数系数的唯一性给出 \(m\) 和 \(u\) 的唯一性。
\end{proof}

\begin{theorem}[恒等定理]
	\label{g4:thm:0101:identity}
	设 \(\Omega\subset\mathbb C\) 是区域，\(f,g\in\mathcal O(\Omega)\)。若集合 \(\{z:f(z)=g(z)\}\) 在 \(\Omega\) 内有聚点，则 \(f=g\) 在整个 \(\Omega\) 上成立。特别地，非恒零全纯函数在区域中的零点离散。
\end{theorem}
\begin{proof}
	令 \(h=f-g\)。由\cref{g4:prop:0101:zero-order}，零点的聚点附近必有 \(h=0\)。令 \(E\) 为所有具有恒零邻域的点的集合，则 \(E\) 非空且开。若 \(p\in\overline E\cap\Omega\)，且 \(p\notin E\)，任意靠近 \(p\) 的 \(E\) 中点都是 \(h\) 的零点，因而 \(p\) 是零点聚点；前述局部结论又给出 \(p\in E\)，矛盾。所以 \(E\) 在 \(\Omega\) 中也闭。连通性给出 \(E=\Omega\)。
\end{proof}

连通条件用于把一处的相等传播到整个区域；内聚点条件则用于启动传播。函数 \(\sin(\pi/z)\) 在穿孔圆盘上有零点 \(1/n\) 向缺失的原点聚集，却并非恒零。由幂级数定义的正弦全纯，故这个例子确实只在聚点的位置上不同。

\ProseParagraphBreak

为了把 \((z-a)^m u(z)\) 改写为新坐标的 \(m\) 次幂，需要先说明一个导数非零的候选坐标为何可逆。以下证明也给出逆函数的全纯性。

\begin{theorem}[全纯逆函数定理]
	\label{g4:thm:0101:inverse}
	设 \(f\) 在 \(a\in\mathbb C\) 的邻域全纯，且 \(f'(a)\ne0\)。则存在 \(a\) 和 \(f(a)\) 的开邻域 \(V,W\)，使 \(f:V\to W\) 双射，其逆 \(g\) 全纯，且
	\[
		g'(w)=\frac1{f'(g(w))}.
	\]
\end{theorem}
\begin{proof}
	通过平移和除以 \(f'(a)\)，化为 \(a=f(a)=0\)、\(f'(0)=1\)，并写 \(f(z)=z+h(z)\)。由已证的导数连续性，取 \(r>0\)，使 \(h(0)=0\) 且 \(|h'|\le q<1\) 在 \(\overline{D(0,r)}\) 上成立。沿直线段积分得到
	\[
		|h(z)-h(\zeta)|\le q|z-\zeta|,\qquad
		|f(z)-f(\zeta)|\ge(1-q)|z-\zeta|.
	\]
	故 \(f\) 在此圆盘上单射。对 \(|w|<(1-q)r\)，迭代 \(z_{n+1}=w-h(z_n)\)，从 \(z_0=0\) 开始。映射把闭圆盘映入自身，且相邻差按因子 \(q\) 缩小。几何级数估计说明 \((z_n)\) 是 Cauchy 列，其极限 \(z\) 满足 \(f(z)=w\)；又 \(|z|\le |w|/(1-q)<r\)。同一压缩估计保证解唯一。

	令 \(W=D(0,(1-q)r)\)、\(V=f^{-1}(W)\cap D(0,r)\)，便得到所需双射。上述下界使逆 \(g\) 为 Lipschitz 连续。若 \(w'\to w\)，\(z'=g(w')\)、\(z=g(w)\)，则
	\[
		\frac{g(w')-g(w)}{w'-w}
		=\left(\frac{f(z')-f(z)}{z'-z}\right)^{-1}
		\longrightarrow\frac1{f'(z)}.
	\]
	这里 \(|f'(z)|\ge1-q\)，所以极限合法。还原归一化即得结论。
\end{proof}

\begin{theorem}[局部幂模型]
	\label{g4:thm:0101:local-power}
	设 \(\Omega\subset\mathbb C\) 是区域，\(f\in\mathcal O(\Omega)\) 非恒定，\(a\in\Omega\)。令 \(m=\operatorname{ord}_a(f-f(a))\)。则存在 \(a\) 附近的全纯坐标 \(w=\phi(z)\)，满足 \(\phi(a)=0\)、\(\phi'(a)\ne0\)，并且
	\begin{equation}
		\label{g4:eq:0101:local-power}
		f(z)=f(a)+\phi(z)^m.
	\end{equation}
	在充分小的坐标邻域 \(V=\phi^{-1}(D(0,\rho))\) 内，每个 \(0<|b-f(a)|<\rho^m\) 恰有 \(m\) 个不同原像；\(f(a)\) 的唯一原像为 \(a\)，重数为 \(m\)。
\end{theorem}
\begin{proof}
	恒等定理保证 \(f-f(a)\) 不会在 \(a\) 附近恒零。由零点分解，写成 \((z-a)^m u(z)\)。缩小圆盘使 \(|u(z)/u(a)-1|<1\)。级数
	\[
		L(\xi)=\sum_{n=1}^\infty\frac{(-1)^{n+1}}n\xi^n\qquad(|\xi|<1)
	\]
	满足 \(L'=1/(1+\xi)\)、\(L(0)=0\)。对 \(e^{L(\xi)}/(1+\xi)\) 求导得到零，由初值可知 \(e^{L(\xi)}=1+\xi\)。任取 \(c^m=u(a)\)，令
	\[
		v(z)=c\exp\left(\frac1m L\left(\frac{u(z)}{u(a)}-1\right)\right),\qquad
		\phi(z)=(z-a)v(z).
	\]
	便有 \(v^m=u\)、\(\phi'(a)=c\ne0\)。\cref{g4:thm:0101:inverse}使 \(\phi\) 在缩小邻域上成为全纯可逆坐标，且给出\eqref{g4:eq:0101:local-power}。再把坐标像缩为 \(D(0,\rho)\)。方程 \(w^m=b-f(a)\) 的 \(m\) 个根都在这个圆盘内且互不相同，映回便得原像计数。
\end{proof}

\begin{theorem}[开映射定理]
	\label{g4:thm:0101:open-mapping}
	区域 \(\Omega\subset\mathbb C\) 上的非恒定全纯函数把开集映为开集。它在点 \(a\) 附近单射，当且仅当 \(f'(a)\ne0\)。
\end{theorem}
\begin{proof}
	对任意开集中的点 \(a\)，在其中取\cref{g4:thm:0101:local-power}的坐标邻域。其像包含 \(f(a)\) 的圆盘，故像开。若 \(f'(a)=0\)，局部幂模型中的 \(m\ge2\)，足够靠近 \(f(a)\) 的非中心点在任意足够小邻域内有多个原像，不能局部单射。若 \(f'(a)\ne0\)，逆函数定理给出局部单射。
\end{proof}

对非恒定全纯函数，\(f'(a)=0\) 的点称为\term[linjiedian]{临界点}[critical point]，其像称为临界值；局部幂模型中的整数 \(m\) 记录这里的重数。譬如 \(f(z)=z^3-3z\) 在 \(a=1\) 处满足
\[
	f(z)+2=(z-1)^2(z+2),
\]
故局部重数为 \(2\)；在原点导数为 \(-3\)，局部重数为 \(1\)。临界点附近的二重性与函数作为三次多项式的整体次数是不同的计数。

\subsection{全纯逆函数及分支}
\label{g4:subsec:0101:04}
逆函数定理给出的逆只在一个邻域内有效。当目标点移动时，这些局部逆可能接续，也可能绕一圈后互换；这正是下一节解析延拓所要记录的现象。

\begin{definition}
	\label{g4:def:0101:biholomorphic}
	设 \(U,V\subset\mathbb C\) 为开集。若 \(f:U\to V\) 是全纯双射且逆映射也全纯，则称 \(f\) 为\term[shuangquanchunyingshe]{双全纯映射}[biholomorphic map]。区域上的单射全纯函数称为\term[danyehanshu]{单叶函数}[univalent function]。若 \(f\) 全纯且 \(f'\) 处处非零，则它局部保有向角，称为保向\term[gongxingyingshe]{共形映射}[conformal map]；本章讨论的共形映照在声称区域等价时还要求双射。
\end{definition}

由\cref{g4:thm:0101:open-mapping}，单叶函数的导数处处非零，像为开集；全局集合逆在每点附近都等于\cref{g4:thm:0101:inverse}给出的全纯局部逆，所以单叶函数总是到其像的双全纯映射。保角性则由\eqref{g4:eq:0101:complex-linear}直接得到：若两条正则曲线在 \(a\) 的切向量为 \(v_1,v_2\ne0\)，两条像曲线的切向量为 \(f'(a)v_1,f'(a)v_2\)，其比值仍为 \(v_2/v_1\)，故有向夹角不变。与此同时，长度元乘以 \(|f'(a)|\)，面积元乘以 \(|f'(a)|^2\)。

\begin{example}
	\label{g4:exm:0101:square-branches}
	设 \(S=\mathbb C\setminus(-\infty,0]\)。每个 \(z\in S\) 唯一写成 \(r e^{\mathrm i\theta}\)，其中 \(r>0\)、\(-\pi<\theta<\pi\)。函数
	\[
		s(z)=\sqrt r\,e^{\mathrm i\theta/2}
	\]
	连续，取值于右半平面，且 \(s(z)^2=z\)。在每个 \(z_0\in S\)，平方映射在 \(s(z_0)\ne0\) 处有全纯局部逆；连续性保证 \(s\) 在充分小邻域内选中的正是这个逆，因此 \(s\) 全纯。另一个根为 \(-s\)。它们是平方根的两个\term[fenzhi]{分支}[branch]，即在指定区域上单值全纯地选择根。
\end{example}

这两个分支不能作为平方根延伸到整个 \(\mathbb C\setminus\{0\}\)。若存在连续平方根 \(s\)，沿 \(z=e^{\mathrm i t}\)、\(0\le t\le2\pi\) 考察
\[
	q(t)=s(e^{\mathrm i t})e^{-\mathrm i t/2}.
\]
它连续且平方为 \(1\)，故在区间上恒为 \(1\) 或恒为 \(-1\)。于是 \(s(e^{2\pi\mathrm i})=-s(1)\)，与起终点相同矛盾。删去一条射线使辐角可以单值选取；保留整个穿孔平面，则沿环路的返回变化无法消去。

\ProseParagraphBreak

导数处处非零也不能保证整体单射。例如指数映射满足 \((e^z)'=e^z\ne0\)，但 \(e^{z+2\pi\mathrm i}=e^z\)。局部幂模型在每点都是一次幂，整体上仍有无穷多个逆分支。要把局部信息用于整体构造，接下来需要两种工具：边界积分对原像的计数，以及沿路径接续函数芽时的唯一性。
% 纲要条目排为编号小节标题，并进入目录与书签；说明段仍保留于纲要注释中。

% !TeX root = ../../Book4.tex
% 纲要标题：### 1.2 留数、次数与解析延拓
% 纲要位置：几何学计划纲要.md，第 3488 行。
% 正文状态：初稿；尚未正式审核。

% 纲要内容开始
% 
% * 亚纯函数和主部
% * 留数定理及辐角原理
% * Rouché 与映射次数
% * 函数芽的延拓和路径依赖
% 
% 纲要内容结束

\section{留数、次数与解析延拓}
\label{g4:sec:0102}

围绕一个零点走一周时，函数值可能绕原点走若干周；沿一条闭路延拓局部定义的对数时，回到起点的值又可能增加 \(2\pi\mathrm{i}\)。这两种现象都把局部解析数据变成了沿闭路测得的整体信息。前者由对数导数的积分计数，后者则要求比较不同路径上的函数芽。本节先把孤立奇点写成可积分的级数，再讨论零点计数和解析延拓。

\subsection{亚纯函数和主部}
\label{g4:subsec:0102:01}
在圆盘内，全纯函数由非负整数次幂展开；若中心被挖去，\(1/(z-a)\) 已经表明必须允许负次幂。负次幂记录了函数在缺失点附近无法由普通幂级数描述的部分。先在整个环域中构造这种展开，再令内半径趋于零，便能同时处理极点和更复杂的孤立奇点。

\begin{theorem}[Laurent 展开]\label{g4:thm:0102:laurent}
	设 \(0\leq r<R\leq\infty\)，函数 \(f\) 在环域
	\[
		A=\{z\in\mathbb C:r<|z-a|<R\}
	\]
	上全纯。则存在唯一的一列复数 \((c_n)_{n\in\mathbb Z}\)，使得
	\[
		f(z)=\sum_{n=-\infty}^{\infty}c_n(z-a)^n,
		\qquad
		c_n=\frac{1}{2\pi\mathrm{i}}
		\int_{|\zeta-a|=\rho}\frac{f(\zeta)}{(\zeta-a)^{n+1}}\,\mathrm{d}\zeta
	\]
	对任意 \(r<\rho<R\) 成立，积分圆周取逆时针方向。正次与负次两部分分别在 \(A\) 的每个紧子集上一致绝对收敛。
\end{theorem}
这个展开称为\term[laurentzhankai]{Laurent 展开}[Laurent expansion]。\footnote{Pierre Alphonse Laurent（1813—1854），法国数学家、工程师，研究复函数的级数展开与变分问题。}
\begin{proof}
	关键是把环域的外边界和内边界分开：外圆贡献非负次幂，内圆贡献负次幂。固定 \(z\in A\)，选取
	\[
		r<\rho_1<|z-a|<\rho_2<R.
	\]
	在两圆之间的闭环域上应用\cref{g4:prop:0101:cauchy-domain}，内边界的正向是顺时针方向，故
	\[
		f(z)=\frac{1}{2\pi\mathrm{i}}
		\int_{|\zeta-a|=\rho_2}\frac{f(\zeta)}{\zeta-z}\,\mathrm{d}\zeta
		-\frac{1}{2\pi\mathrm{i}}
		\int_{|\zeta-a|=\rho_1}\frac{f(\zeta)}{\zeta-z}\,\mathrm{d}\zeta.
	\]
	外圆上用几何级数
	\[
		\frac{1}{\zeta-z}
		=\sum_{n=0}^{\infty}\frac{(z-a)^n}{(\zeta-a)^{n+1}},
	\]
	内圆上则用
	\[
		\frac{1}{\zeta-z}
		=-\sum_{k=0}^{\infty}\frac{(\zeta-a)^k}{(z-a)^{k+1}}.
	\]
	在任意较小的闭环域 \(\rho_1<r'\leq|z-a|\leq R'<\rho_2\) 上，两式的公比分别不超过 \(R'/\rho_2\) 和 \(\rho_1/r'\)，均严格小于 \(1\)。因此可以逐项积分，所得两级数在该闭环域上一致绝对收敛。

	对每个整数 \(n\)，被积函数 \(f(\zeta)/(\zeta-a)^{n+1}\) 在 \(A\) 上全纯。在两条半径之间的闭环域上应用边界积分为零的结论，说明其圆周积分与半径无关，故上述系数都可用同一个 \(\rho\) 表示。任意紧子集均包含在一个这样的较小闭环域中，得到所需收敛性。若已有另一展开，在 \(|z-a|=\rho\) 上逐项乘以 \((z-a)^{-m-1}\) 积分，只有第 \(m\) 项的积分非零，因而必恢复同一系数 \(c_m\)，唯一性得证。
\end{proof}

取 \(r=0\) 时，点 \(a\) 称为 \(f\) 的\term[guliqidian]{孤立奇点}[isolated singularity]。展开中的
\[
	\sum_{k=1}^{\infty}c_{-k}(z-a)^{-k}
\]
称为该奇点的\term[zhubu]{主部}[principal part]。若主部为零，非负次幂部分给出穿过 \(a\) 的全纯延拓，称 \(a\) 为\term[kequqidian]{可去奇点}[removable singularity]。若主部恰有有限个非零项，且最负的非零项为 \(c_{-m}(z-a)^{-m}\)，则称 \(a\) 为 \(m\) 阶\term[jidian]{极点}[pole]；等价地，
\[
	f(z)=(z-a)^{-m}u(z),\qquad u\text{ 在 }a\text{ 附近全纯},\quad u(a)\neq0.
\]
主部有无穷多个非零项时，称 \(a\) 为\term[benzhiqidian]{本性奇点}[essential singularity]。这三种情况穷尽孤立奇点，因为它们正是负次幂系数的三种可能。

\begin{proposition}\label{g4:prop:0102:singularity-tests}
	设 \(f\) 在 \(0<|z-a|<R\) 上全纯。则 \(a\) 可去当且仅当 \(f\) 在某个去心邻域上有界；\(a\) 是极点当且仅当 \(|f(z)|\rightarrow\infty\) 当 \(z\rightarrow a\)。
\end{proposition}
\begin{proof}
	若 \(|f|\leq M\)，则对于每个 \(k\geq1\)，Laurent 系数公式给出
	\[
		|c_{-k}|\leq M\rho^k.
	\]
	令 \(\rho\rightarrow0\)，所有负次幂系数都为零。反向由全纯延拓的连续性得到。极点的局部表示直接给出模趋于无穷。反之，若 \(|f(z)|\rightarrow\infty\)，则 \(f\) 在充分小的去心圆盘上无零点，且 \(1/f\rightarrow0\)。可去性判据把 \(1/f\) 延拓为在 \(a\) 取零的全纯函数；它在去心圆盘上不为零，故由\cref{g4:thm:0101:identity}，该零点具有有限阶。取倒数便得到极点表示。
\end{proof}

\begin{definition}\label{g4:def:0102:meromorphic}
	设 \(\Omega\subset\mathbb C\) 为域。若 \(P\subset\Omega\) 在 \(\Omega\) 内无聚点，\(f\) 在 \(\Omega\setminus P\) 上全纯，且每个 \(a\in P\) 都是 \(f\) 的极点，则称 \(f\) 为 \(\Omega\) 上的\term[yachunhanshu]{亚纯函数}[meromorphic function]，在极点处约定 \(f(a)=\infty\)。允许 \(P=\varnothing\)，所以全纯函数也是亚纯函数。若 \(f\) 不恒为零，则在每个 \(a\in\Omega\) 附近可以唯一地写出阶数
	\[
		f(z)=(z-a)^m u(z),\qquad m\in\mathbb Z,\quad u(a)\neq0,
	\]
	其中 \(u\) 全纯。整数 \(m\) 记为 \(\operatorname{ord}_a f\)：零点处为正，极点处为负，其余点为零。\symidx{\(\operatorname{ord}_a f\)：亚纯函数的阶数}
\end{definition}

若 \(g,h\) 在域上全纯且 \(h\) 不恒为零，则 \(g/h\) 在 \(h\neq0\) 处定义，在共同零点处约去相同次数的因子后，所得函数是亚纯的。局部阶数满足
\[
	\operatorname{ord}_a(fg)=\operatorname{ord}_a f+\operatorname{ord}_a g,
	\qquad
	\operatorname{ord}_a(f/g)=\operatorname{ord}_a f-\operatorname{ord}_a g
\]
（各函数均不恒为零）。这些等式把乘除法转成整数的加减，是后续用除子记录零极点的基础。

\begin{example}\label{g4:exm:0102:laurent-domains}
	同一函数在不同环域上的 Laurent 展开可以不同。例如
	\[
		\frac{1}{z(z-1)}=-\frac1z+\frac1{z-1}
		=\begin{cases}
			-\dfrac1z-\displaystyle\sum_{n=0}^{\infty}z^n,&0<|z|<1,\\[6pt]
			\displaystyle\sum_{n=2}^{\infty}z^{-n},&|z|>1.
		\end{cases}
	\]
	第二式不能用于读取 \(0\) 处的主部，因为它的环域并不贴近 \(0\)。另一方面，\(e^{1/z}=\sum_{n=0}^{\infty}z^{-n}/n!\) 在 \(0\) 处具有本性奇点，所以它不是整个复平面上的亚纯函数；在去心平面上它仍是全纯函数。奇点的类型必须连同所讨论的定义域一起判断。
\end{example}

\subsection{留数定理及辐角原理}
\label{g4:subsec:0102:02}
对圆周积分，Laurent 级数中只有 \((z-a)^{-1}\) 一项留下贡献。因此，一个系数便能控制围绕奇点的积分。

\begin{definition}\label{g4:def:0102:residue}
	设 \(f\) 在 \(a\) 的去心邻域上全纯。其 Laurent 展开的系数 \(c_{-1}\) 称为 \(f\) 在 \(a\) 处的\term[liushu]{留数}[residue]，记为 \(\operatorname{Res}_{z=a}f\)。\symidx{\(\operatorname{Res}_{z=a}f\)：留数}对于充分小的正向圆周，
	\[
		\operatorname{Res}_{z=a}f
		=\frac{1}{2\pi\mathrm{i}}\int_{|z-a|=\rho}f(z)\,\mathrm{d}z.
	\]
\end{definition}

若 \(a\) 是至多 \(m\) 阶极点，则把 \((z-a)^m f(z)\) 展开成普通幂级数，可得实用公式
\[
	\operatorname{Res}_{z=a}f
	=\frac{1}{(m-1)!}
	\left.\frac{\mathrm{d}^{m-1}}{\mathrm{d}z^{m-1}}
	\bigl((z-a)^m f(z)\bigr)\right|_{z=a}.
\]
这里先把括号内函数全纯延拓到 \(a\)，再求导。特别地，若 \(g,h\) 在 \(a\) 附近全纯，\(h(a)=0\)、\(h'(a)\neq0\)，则
\[
	\operatorname{Res}_{z=a}\frac{g(z)}{h(z)}=\frac{g(a)}{h'(a)}.
\]

\begin{theorem}[留数定理]\label{g4:thm:0102:residue}
	设 \(G\subset\mathbb C\) 是有界域，边界由有限条互不相交的分片 \(C^2\) 正则简单闭曲线组成，角点内角在 \((0,2\pi)\) 内，并按使 \(G\) 位于左侧的方向定向。设 \(f\) 在 \(\overline G\) 的一个开邻域中，除有限个内点 \(a_1,\ldots,a_N\in G\) 外全纯，各 \(a_j\) 都是孤立奇点。则
	\[
		\int_{\partial G}f(z)\,\mathrm{d}z
		=2\pi\mathrm{i}\sum_{j=1}^N\operatorname{Res}_{z=a_j}f.
	\]
\end{theorem}
这就是\term[liushudingli]{留数定理}[residue theorem]的有限多边界版本；内边界取顺时针方向。这里每条边界弧的正则性指其 \(C^2\) 参数化的一阶导数不为零，角点内角指两条相邻边界弧在区域一侧所夹的角；这些条件与\eqref{g4:eq:0101:green}采用的边界条件一致。
\begin{proof}
	在每个 \(a_j\) 周围挖去一个闭圆盘，令这些圆盘两两不交且包含于 \(G\)。剩余闭区域的一个邻域上 \(f\) 全纯，由\cref{g4:prop:0101:cauchy-domain}的边界积分结论，其总边界积分为零。因此
	\[
		\int_{\partial G}f(z)\,\mathrm{d}z
		=\sum_{j=1}^N\int_{|z-a_j|=\rho_j}f(z)\,\mathrm{d}z,
	\]
	右侧各圆均取逆时针方向。每个圆周积分由 Laurent 系数公式恰等于对应留数的 \(2\pi\mathrm{i}\) 倍。若挖孔后出现多个连通分支，对各分支相加得到同一等式。
\end{proof}

留数也可以内在地赋给亚纯一形式 \(f(z)\,\mathrm{d}z\)。若换用局部全纯坐标 \(z=a+\phi(w)\)，其中 \(\phi(0)=0\)、\(\phi'(0)\neq0\)，则
\[
	\operatorname{Res}_{w=0}\bigl(f(a+\phi(w))\phi'(w)\bigr)
	=\operatorname{Res}_{z=a}f(z).
\]
为看清这一点，把亚纯函数的有限主部逐项处理。除 \((z-a)^{-1}\mathrm{d}z\) 外，各负次幂项都是某个亚纯函数的微分，换坐标后仍是微分，其 Laurent 展开的 \(w^{-1}\mathrm{d}w\) 系数为零；全纯部分也不产生负次幂。剩下的一项满足
\[
	\frac{\phi'(w)}{\phi(w)}\,\mathrm{d}w
	=\frac{\mathrm{d}w}{w}+\frac{v'(w)}{v(w)}\,\mathrm{d}w,
	\qquad \phi(w)=wv(w),\quad v(0)\neq0,
\]
故留数为 \(1\)。这一不变性解释了为什么曲面上的留数应当属于一形式：只改变函数表达式而遗漏微分因子，数值一般会改变。

\ProseParagraphBreak

接下来把积分用于数零点。对一条避开 \(0\) 的分片 \(C^1\) 闭曲线 \(\eta:[0,1]\rightarrow\mathbb C\setminus\{0\}\)，记
\[
	\operatorname{Ind}(\eta,0)
	\coloneqq\frac{1}{2\pi\mathrm{i}}\int_\eta\frac{\mathrm{d}w}{w}.
\]
这个数称为曲线关于 \(0\) 的\term[raoshu]{绕数}[winding number]。它是整数：把参数区间分成有限个小段，使每段的像落在某个可选连续辐角的小圆盘中，在相邻段上调整辐角的 \(2\pi\) 整数倍，使之连续拼接。于是有分片 \(C^1\) 的函数 \(r(t)>0\)、\(\theta(t)\in\mathbb R\)，满足 \(\eta(t)=r(t)e^{\mathrm{i}\theta(t)}\)，且
\[
	\int_\eta\frac{\mathrm{d}w}{w}
	=\log r(1)-\log r(0)+\mathrm{i}\bigl(\theta(1)-\theta(0)\bigr)
	=2\pi\mathrm{i}k,
\]
其中 \(k\in\mathbb Z\)，因为 \(\eta(1)=\eta(0)\)。所以绕数记录的是总辐角增量，而不是曲线图形的某种单调转动次数。

\begin{theorem}[辐角原理]\label{g4:thm:0102:argument}
	设 \(G\subset\mathbb C\) 是有界域，边界由有限条互不相交的分片 \(C^2\) 正则简单闭曲线组成，角点内角在 \((0,2\pi)\) 内，并取正向。设 \(f\) 在包含 \(\overline G\) 的某个域上亚纯且不恒为零，在 \(\partial G\) 上没有零点或极点。则
	\[
		\frac{1}{2\pi\mathrm{i}}\int_{\partial G}\frac{f'(z)}{f(z)}\,\mathrm{d}z
		=\sum_{a\in G}\operatorname{ord}_a f
		=N_G(f)-P_G(f),
	\]
	其中 \(N_G(f)\)、\(P_G(f)\) 分别按重数统计 \(G\) 内的零点数和极点数。左侧也等于各条边界像曲线 \(f\circ\gamma\) 关于 \(0\) 的绕数之和。
\end{theorem}
这里的\term[fujiaoyuanli]{辐角原理}[argument principle]把解析阶数变为整数值的边界数据。
\begin{proof}
	零点和极点在 \(\overline G\) 的邻域中离散，而边界上没有这两类点；紧性保证 \(G\) 内只有有限个。若
	\[
		f(z)=(z-a)^m u(z),\qquad u(a)\neq0,
	\]
	则
	\[
		\frac{f'(z)}{f(z)}=\frac{m}{z-a}+\frac{u'(z)}{u(z)},
	\]
	后项在 \(a\) 附近全纯。因此 \(f'/f\) 在 \(a\) 的留数为 \(m\)，即零点阶数取正、极点阶数取负。对 \(f'/f\) 应用\cref{g4:thm:0102:residue}便得等式。沿每条边界曲线作代换 \(w=f(z)\)，即得到像曲线的绕数表达式。
\end{proof}

\subsection{Rouché 与映射次数}
\label{g4:subsec:0102:03}
零点的位置可能随函数扰动而移动，而绕数只取整数值。只要扰动过程中边界像不经过原点，整数就无法连续改变。这给出了不必求出零点位置的计数方法。

\begin{theorem}[Rouché 定理]\label{g4:thm:0102:rouche}
	设 \(G\subset\mathbb C\) 是有界域，边界由有限条互不相交的分片 \(C^2\) 正则简单闭曲线组成，角点内角在 \((0,2\pi)\) 内。函数 \(f,g\) 在 \(\overline G\) 的一个开邻域上全纯，并且
	\[
		|g(z)-f(z)|<|f(z)|\qquad(z\in\partial G).
	\]
	则 \(f\) 与 \(g\) 在 \(G\) 内的零点总数相同，零点均按重数统计。
\end{theorem}
这一零点比较结论称为\term[rouchedingli]{Rouché 定理}[Rouché's theorem]。\footnote{Eugène Rouché（1832—1910），法国数学家，研究复函数的零点计数，并参与编写几何学教材。}
\begin{proof}
	以线段插值连接两个函数：\(f_t=f+t(g-f)\)，\(0\leq t\leq1\)。边界上有
	\[
		|f_t(z)|\geq|f(z)|-t|g(z)-f(z)|>0.
	\]
	边界紧，且 \((t,z)\mapsto f_t(z)\) 连续，故分母在 \([0,1]\times\partial G\) 上与零保持正距离。于是
	\[
		n(t)=\frac{1}{2\pi\mathrm{i}}\int_{\partial G}\frac{f_t'(z)}{f_t(z)}\,\mathrm{d}z
	\]
	是 \(t\) 的连续函数。由\cref{g4:thm:0102:argument}，它等于 \(f_t\) 的零点数，因而取整数值。连续的整数值函数在区间上为常数，所以 \(n(0)=n(1)\)。
\end{proof}

\begin{example}\label{g4:exm:0102:rouche-polynomial}
	令 \(p(z)=z^5+3z+1\)。在 \(|z|=1\) 上，
	\[
		|p(z)-3z|=|z^5+1|\leq2<3=|3z|,
	\]
	故单位圆盘内恰有一个零点。在 \(|z|=2\) 上，
	\[
		|p(z)-z^5|=|3z+1|\leq7<32=|z^5|,
	\]
	故半径 \(2\) 的圆盘内恰有五个零点。因此环域 \(1<|z|<2\) 内共有四个零点，均按重数计；两个边界圆上都没有零点。比较项由所选边界上的大小决定，同一个多项式在两条圆周上适合比较的主项并不相同。
\end{example}

设 \(G\) 为边界满足\cref{g4:thm:0102:rouche}条件的有界域，\(f\) 在 \(\overline G\) 的邻域上全纯，且在 \(G\) 上不恒定。对任意 \(w\notin f(\partial G)\)，定义
\begin{equation}\label{g4:eq:0102:degree}
	\deg(f,G,w)\coloneqq
	\frac{1}{2\pi\mathrm{i}}\int_{\partial G}\frac{f'(z)}{f(z)-w}\,\mathrm{d}z
	=\sum_{\substack{z\in G\\f(z)=w}}\operatorname{ord}_z(f-w).
\end{equation}
这是 \(f\) 在 \(G\) 中关于目标点 \(w\) 的\term[yingshecishu]{映射次数}[mapping degree]。等式由辐角原理应用于 \(f-w\) 得到，也说明该数是非负整数。由于积分依 \(w\) 连续，它在 \(\mathbb C\setminus f(\partial G)\) 的每个连通分支上恒定。因而，只要目标点不越过边界像，纤维的重数总和就保持不变。

\ProseParagraphBreak

若 \(w\) 的所有原像都是正则点，即 \(f'(z)\neq0\)，\cref{g4:thm:0101:inverse}给出每个原像附近的一张全纯逆分支，\eqref{g4:eq:0102:degree}便是这些原像的个数。在临界点附近，\cref{g4:thm:0101:local-power}把映射化成 \(\zeta\mapsto\zeta^m\)：非零的小目标值有 \(m\) 个互异原像，目标值到达临界值后它们合并成一个 \(m\) 重原像。次数保存的正是合并前后的这个总数。

\begin{proposition}\label{g4:prop:0102:local-count}
	设 \(f\) 在 \(a\in\mathbb C\) 的邻域上全纯，且 \(f-f(a)\) 在 \(a\) 有 \(m\geq1\) 阶零点。则存在 \(r>0\) 和 \(\varepsilon>0\)，使得对每个满足 \(|w-f(a)|<\varepsilon\) 的 \(w\)，方程 \(f(z)=w\) 在 \(|z-a|<r\) 内恰有 \(m\) 个解，按重数统计。
\end{proposition}
\begin{proof}
	取充分小的 \(r\)，使 \(f-f(a)\) 在 \(0<|z-a|\leq r\) 上无零点，并在闭圆盘的邻域上全纯。令
	\[
		\varepsilon=\min_{|z-a|=r}|f(z)-f(a)|>0.
	\]
	在该圆周上，\(|(f-w)-(f-f(a))|=|w-f(a)|<|f-f(a)|\)。由\cref{g4:thm:0102:rouche}，\(f-w\) 与 \(f-f(a)\) 的零点总数同为 \(m\)。
\end{proof}

这个证明还显示，小扰动保持的是总重数，未必保持互异零点数。例如 \(z^2-t\) 在 \(t\neq0\) 时有两个单零点，在 \(t=0\) 时只有一个二重零点。

\subsection{函数芽的延拓和路径依赖}
\label{g4:subsec:0102:04}
幂级数只在某个圆盘中给出函数，但圆盘的边界未必是函数真正的障碍。例如 \(\sum_{n\geq0}z^n\) 的收敛半径为 \(1\)，它的和 \(1/(1-z)\) 却可以穿过单位圆上除 \(1\) 外的所有点。延拓要保留的是邻域中函数的全部局部信息，而不是某个圆盘或某个展开式。

\begin{definition}\label{g4:def:0102:germ-continuation}
	设 \(a\in\Omega\subset\mathbb C\)，其中 \(\Omega\) 为域。两个在 \(a\) 的邻域上定义的全纯函数，若在某个更小邻域中相等，则视为等价；等价类称为 \(a\) 处的\term[hanshuya]{全纯函数芽}[germ of a holomorphic function]，记为 \([f]_a\)。

	设 \(\gamma:[0,1]\rightarrow\Omega\) 为从 \(a\) 出发的连续路径。一个沿 \(\gamma\) 的\term[jiexiyantuo]{解析延拓}[analytic continuation]由有限分割
	\[
		0=t_0<t_1<\cdots<t_N=1
	\]
	以及圆盘 \(U_j\subset\Omega\) 上的全纯函数 \(f_j\) 给出，满足 \(\gamma([t_{j-1},t_j])\subset U_j\)，第一项代表指定的初始函数芽，且
	\[
		[f_j]_{\gamma(t_j)}=[f_{j+1}]_{\gamma(t_j)}
		\qquad(1\leq j<N).
	\]
	末项给出的 \([f_N]_{\gamma(1)}\) 称为延拓的终点函数芽。
\end{definition}

这里的相等要求是在交点的一个邻域中相等；只要求数值相等无法传递导数和幂级数。两个圆盘的交集是连通的，所以由\cref{g4:thm:0101:identity}，相邻代表实际上在整个交集上一致。分割与圆盘只是记录延拓的工具，可以继续细分或缩小而不改变沿路径所得的函数芽。

\begin{lemma}\label{g4:lem:0102:continuation-unique}
	设 \(\Omega\subset\mathbb C\) 为域，\(\gamma:[0,1]\rightarrow\Omega\) 为连续路径。指定 \(\gamma(0)\) 处的全纯函数芽后，若沿 \(\gamma\) 的解析延拓存在，则每个 \(\gamma(t)\) 处随路径传递的函数芽都唯一，特别是终点函数芽唯一。
\end{lemma}
\begin{proof}
	比较两条延拓链，先对它们的分割取公共细分。在每个小参数区间内，两条链分别由两个圆盘上的函数代表，且路径段同时落在这两个圆盘中。开始时两个代表的芽相同；圆盘交集连通，恒等定理使两个代表在交集上相同，于是这一段的终点芽也相同。接着利用各自延拓链在接点处的相容性，进入下一段。有限次归纳便得到所有分割点的芽相同；把任意给定 \(t\) 加入分割，便知该处的芽也相同。
\end{proof}

固定路径上的唯一性仍允许不同路径给出不同结果。在 \(1\) 附近取满足 \(L(1)=0\)、\(e^{L(z)}=z\) 的局部对数，沿单位圆 \(\gamma(t)=e^{2\pi\mathrm{i}t}\) 延拓。连续辐角在这条路径上是 \(2\pi t\)，故沿途的值为 \(2\pi\mathrm{i}t\)，终点对数芽在 \(1\) 的值变成 \(2\pi\mathrm{i}\)。从同一起点走常值路径，终点值仍为 \(0\)。同样，从 \(\sqrt1=1\) 出发的平方根沿此圆周的值为 \(e^{\pi\mathrm{i}t}\)，一周后变成 \(-1\)。

\ProseParagraphBreak

上述值确实来自解析延拓链。对任意 \(b\neq0\)，在 \(|z-b|<|b|\) 内写成
\[
	L_b(z)=\ell_b+\sum_{n=1}^{\infty}\frac{(-1)^{n+1}}n
	\left(\frac{z-b}{b}\right)^n,\qquad e^{\ell_b}=b.
\]
逐项求导得 \(L_b'(z)=1/z\)，且 \(e^{L_b(z)}/z\) 的导数为零，在 \(b\) 取值 \(1\)，所以 \(e^{L_b(z)}=z\)。相邻圆盘上两个这样的对数相差 \(2\pi\mathrm{i}\) 的整数倍；调整常数后即可沿路径拼接。平方根由 \(e^{L_b/2}\) 得到。这也说明导致路径依赖的是绕过缺失点后的常数变化。

\ProseParagraphBreak

回顾这里的拓扑条件：两条同起终点路径\term[dingduandiantonglun]{定端点同伦}[homotopic relative to endpoints]，是指存在连续变形 \(H(t,s)\)，在变形过程中始点和终点始终固定。域内每条闭路都能这样收缩到常值路径时，称该域\term[danliantong]{单连通}[simply connected]。等价地，域内任意两条同起终点路径都定端点同伦；把一条路径接上另一条的反向路径，便得到这两个表述之间的对应。

\begin{theorem}[单值定理]\label{g4:thm:0102:monodromy}
	设 \(\Omega\subset\mathbb C\) 为域，\(a\in\Omega\)，且 \(\mathfrak f\) 是 \(a\) 处的全纯函数芽。假设 \(\mathfrak f\) 可以沿 \(\Omega\) 中每条从 \(a\) 出发的连续路径解析延拓。则沿两条定端点同伦路径得到的终点函数芽相同。特别地，若 \(\Omega\) 单连通，则存在唯一的全纯函数 \(F\) 定义在整个 \(\Omega\) 上，使 \([F]_a=\mathfrak f\)，且所有上述延拓均给出 \(F\) 的相应局部函数芽。
\end{theorem}
这里的\term[danzhidingli]{单值定理}[monodromy theorem]要求沿每条路径都能延拓。单连通是消除路径选择的拓扑条件，延拓存在性则是关于所给函数芽的解析条件。
\begin{proof}
	先证明一条有限延拓链在路径的小变动下仍然有效。固定从 \(a\) 到 \(b\) 的路径 \(\gamma\)，取延拓链 \((U_j,f_j)\) 和相应分割。每个紧集 \(\gamma([t_{j-1},t_j])\) 都包含于开圆盘 \(U_j\)，故存在 \(\delta_j>0\)，使这一紧集的 \(\delta_j\) 邻域仍包含于 \(U_j\)。令 \(\delta\) 为有限个 \(\delta_j\) 的最小值。若另一条路径 \(\widetilde\gamma\) 具有相同端点，并满足
	\[
		\sup_{0\leq t\leq1}|\widetilde\gamma(t)-\gamma(t)|<\delta,
	\]
	则它的每个相应路径段仍位于 \(U_j\)。相邻函数在圆盘交集上相等，所以同一列 \(f_j\) 也给出沿 \(\widetilde\gamma\) 的延拓；起终点固定，两条路径的初始芽与终点芽分别相同。由\cref{g4:lem:0102:continuation-unique}，这就是 \(\widetilde\gamma\) 的唯一终点芽。

	现在设 \(H:[0,1]^2\rightarrow\Omega\) 是定端点同伦，令 \(\gamma_s(t)=H(t,s)\)，其中 \(H(0,s)=a\)、\(H(1,s)=b\)。按假设，每条 \(\gamma_s\) 都有延拓。固定 \(s_0\)，对 \(\gamma_{s_0}\) 使用上一段的有限链；由 \(H\) 在紧方形上一致连续，\(s\) 充分靠近 \(s_0\) 时，\(\gamma_s\) 与 \(\gamma_{s_0}\) 在上述一致距离下小于 \(\delta\)。因此终点芽关于 \(s\) 局部不变。区间 \([0,1]\) 连通，一个局部不变的取值在全区间上必相同：给定一个值，其原像和补集都是开集，且原像非空。所以 \(s=0\) 与 \(s=1\) 的终点芽相同。

	若 \(\Omega\) 单连通，则任意两条具有相同端点的路径定端点同伦。对每个 \(b\in\Omega\)，用任意一条从 \(a\) 到 \(b\) 的路径延拓，并把终点芽的值记为 \(F(b)\)；刚才的结论保证选择无关。为证明 \(F\) 全纯，固定一条到 \(b\) 的路径，令它的末项代表为圆盘 \(U\) 上的函数 \(h\)。在 \(b\) 周围取更小圆盘 \(V\subset U\)。对于 \(c\in V\)，先沿原路径到 \(b\)，再沿线段到 \(c\)，末段使用同一个函数 \(h\)，便得 \(F(c)=h(c)\)。因此 \(F\) 在每点附近都是全纯的，并具有指定初始芽。唯一性由域上的恒等定理得到。
\end{proof}

例如，\(1/(1-z)\) 在 \(0\) 处的函数芽不能沿 \(\mathbb C\) 中终点为 \(1\) 的路径全纯延拓：若终点附近存在全纯代表 \(h\)，在与前一段重合的开集上便有 \((1-z)h(z)=1\)，恒等定理会使此式在 \(1\) 附近成立，在 \(z=1\) 处矛盾。因此，尽管 \(\mathbb C\) 单连通，单值定理的延拓存在性假设仍不满足。

\begin{proposition}\label{g4:prop:0102:logarithm}
	设 \(\Omega\subset\mathbb C\) 是单连通域，\(f\) 在 \(\Omega\) 上全纯且处处非零。则存在全纯函数 \(L\) 使 \(e^L=f\)。指定 \(a\in\Omega\) 及一个满足 \(e^\ell=f(a)\) 的数 \(\ell\) 后，可唯一选定满足 \(L(a)=\ell\) 的 \(L\)。此外，对每个整数 \(m\geq1\)，\(f\) 有全纯 \(m\) 次根；任意两个这样的根相差一个常数 \(m\) 次单位根。
\end{proposition}
\begin{proof}
	每个 \(b\in\Omega\) 附近，都可选圆盘 \(U_b\subset\Omega\)，使 \(|f(z)/f(b)-1|<1\)。选择 \(\ell_b\) 满足 \(e^{\ell_b}=f(b)\)，并令
	\[
		L_b(z)=\ell_b+\sum_{n=1}^{\infty}\frac{(-1)^{n+1}}n
		\left(\frac{f(z)}{f(b)}-1\right)^n.
	\]
	与前面的局部对数计算相同，\(L_b'=f'/f\)、\(e^{L_b}=f\)。若两个局部对数定义在相交圆盘上，其差的指数恒为 \(1\)，故差只取 \(2\pi\mathrm{i}\mathbb Z\) 中的值；连续性和交集连通性说明这个差是常数。

	从 \(a\) 出发任取一条连续路径。用上述圆盘覆盖它的紧像，取参数区间的有限细分，使每个小路径段落在某个圆盘中。先把第一项对数的值调整为指定的 \(\ell\)，再依次加上 \(2\pi\mathrm{i}\) 的整数倍，使相邻对数在交点的芽相同。这样便构成沿该路径的解析延拓。由\cref{g4:thm:0102:monodromy}，这些局部对数粘成整个 \(\Omega\) 上的全纯 \(L\)。恒等式 \(e^L=f\) 在局部成立，故处处成立。两个全纯对数之差在连通域上是常数 \(2\pi\mathrm{i}k\)，指定初值便确定唯一选择。

	令 \(g=e^{L/m}\)，则 \(g^m=f\)。若 \(\widetilde g^m=f\)，因为 \(f\neq0\)，比值 \(\widetilde g/g\) 全纯且只取有限个 \(m\) 次单位根中的值，连通性使之恒定。
\end{proof}

全纯对数的存在还可由闭路积分检测。在任意域 \(\Omega\) 上，无零点全纯函数 \(f\) 有全纯对数，当且仅当
\[
	\int_\gamma\frac{f'(z)}{f(z)}\,\mathrm{d}z=0
	\qquad\text{对每条分片 }C^1\text{ 闭路 }\gamma\subset\Omega.
\]
必要性由 \(f'/f=L'\) 得到。反之，固定 \(a\) 及 \(e^\ell=f(a)\)，定义
\[
	L(z)=\ell+\int_a^z\frac{f'(\zeta)}{f(\zeta)}\,\mathrm{d}\zeta.
\]
域中任意两点可用折线路径连接，闭路积分为零使定义与所选折线路径无关。在小圆盘内沿最后一条线段求导，得 \(L'=f'/f\)，从而 \((fe^{-L})'=0\)，初值使 \(fe^{-L}=1\)。在 \(\mathbb C\setminus\{0\}\) 上取 \(f(z)=z\)，单位圆的积分为 \(2\pi\mathrm{i}\)，所以不存在全域对数。留数、绕数与延拓的路径依赖在这个例子中给出了同一障碍的三种表达。
% 纲要条目排为编号小节标题，并进入目录与书签；说明段仍保留于纲要注释中。

% !TeX root = ../../Book4.tex
% 纲要标题：### 1.3 极值原理和刚性
% 纲要位置：几何学计划纲要.md，第 3495 行。
% 正文状态：初稿；尚未正式审核。

% 纲要内容开始
% 
% * 最大模及调和函数
% * Schwarz 引理与等号
% * Schwarz–Pick 的度量形式
% * 双全纯映射与自同构
% 
% 纲要内容结束

\section{极值原理和刚性}
\label{g4:sec:0103}

一个全纯函数在小圆盘内可以有任意高阶的零点，但它的模不能在内点形成孤立的峰。这个限制把局部解析性转化为整体控制：边界上的估计约束内部，极值处的等号又往往迫使整个函数具有特殊形式。圆盘自映射最清楚地体现了这一点。先固定一个点，把函数与恒等映射比较，再改变所固定的点，就能得到适合所有圆盘自映射的长度尺度。

\subsection{最大模及调和函数}
\label{g4:subsec:0103:01}
先从几何图像理解模的极值。若 \(f\) 非常值，则\cref{g4:thm:0101:open-mapping}说明它把开集映成开集。假如 \(f(a)\ne0\)，一个包含 \(f(a)\) 的小圆盘必有模大于 \(\lvert f(a)\rvert\) 的点；因此像集不可能在 \(f(a)\) 处贴住以原点为中心的圆而完全留在圆内。

\begin{theorem}[最大模原理]
	\label{g4:thm:0103:maximum-modulus}
	设 \(\Omega\subset\mathbb C\) 是区域，\(f\colon\Omega\rightarrow\mathbb C\) 全纯。若 \(\lvert f\rvert\) 在 \(\Omega\) 的某一点取得局部最大值，则 \(f\) 在 \(\Omega\) 上为常值。若进一步假设 \(\Omega\) 有界，且 \(f\) 连续延拓到 \(\overline\Omega\)，则
	\[
		\max_{\overline\Omega}\lvert f\rvert
		=\max_{\partial\Omega}\lvert f\rvert.
	\]
\end{theorem}

\begin{proof}
	设 \(a\in\Omega\)，存在邻域 \(U\subset\Omega\) 使 \(\lvert f(z)\rvert\leq\lvert f(a)\rvert\)。若 \(f\) 非常值，开映射定理使 \(f(U)\) 成为包含 \(f(a)\) 的开集，因而包含模大于 \(\lvert f(a)\rvert\) 的点，矛盾。这也覆盖 \(f(a)=0\) 的情形；此时局部上界本身就迫使 \(f|_U=0\)。

	在有界域的情形，\(\overline\Omega\) 是非空紧集，连续函数 \(\lvert f\rvert\) 在其上取得最大值。若最大值在内点取得，前一结论使 \(f\) 为常值，边界也取得同一值；否则最大值直接在边界取得。有界非空区域的边界非空，所以右端有意义。
\end{proof}

这里的\term[zuidamoyuanli]{最大模原理}[maximum modulus principle]控制的是复函数的模。实部和虚部同样受到极值限制，其原因可以用一个实偏微分方程表达。

\begin{definition}
	\label{g4:def:0103:harmonic}
	设 \(U\subset\mathbb R^2\) 开，\(u\colon U\rightarrow\mathbb R\) 为 \(C^2\) 函数。若
	\[
		\Delta u=-(u_{xx}+u_{yy})=0,
	\]
	则称 \(u\) 为\term[tiaohehanshu]{调和函数}[harmonic function]。这里 \(\Delta\) 沿用全书的非负 Laplace 算子约定。\footnote{Pierre-Simon Laplace（1749--1827），法国数学家，研究天体力学、概率论与位势理论。}若实函数 \(v\) 使 \(u+\mathrm{i}v\) 在 \(U\subset\mathbb C\) 上全纯，则称 \(v\) 为 \(u\) 的\term[tiaohegonge]{调和共轭}[harmonic conjugate]。
\end{definition}

全纯函数 \(f=u+\mathrm{i}v\) 的实部与虚部都是调和函数。事实上，全纯函数由幂级数表示，因而可以连续求任意阶偏导；对 Cauchy–Riemann 方程 \(u_x=v_y\)、\(u_y=-v_x\) 再求导，得到 \(u_{xx}+u_{yy}=v_{yx}-v_{xy}=0\)，对 \(v\) 也有同样结论。反方向的困难在于能否由 \(u\) 恢复 \(v\)：所需导数必须满足
\[
	v_x=-u_y,\qquad v_y=u_x.
\]
调和方程正是这两个导数的相容条件。

\begin{proposition}
	\label{g4:prop:0103:harmonic-conjugate}
	设 \(u\) 是区域 \(\Omega\subset\mathbb C\) 上的实值调和函数。则每一点附近都有调和共轭。在整个 \(\Omega\) 上存在调和共轭，当且仅当对每条分片 \(C^1\) 闭曲线 \(\gamma\subset\Omega\) 都有
	\begin{equation}
		\label{g4:eq:0103:conjugate-period}
		\int_\gamma(-u_y\,\mathrm{d}x+u_x\,\mathrm{d}y)=0.
	\end{equation}
	若全局调和共轭存在，则相差一个实常数。
\end{proposition}

\begin{proof}
	先在包含给定点 \((x_0,y_0)\) 且闭包位于 \(\Omega\) 内的小开矩形上构造。沿水平线段再沿竖直线段积分，令
	\[
		v(x,y)=-\int_{x_0}^{x}u_y(s,y_0)\,\mathrm{d}s
		+\int_{y_0}^{y}u_x(x,t)\,\mathrm{d}t.
	\]
	由微积分基本定理，\(v_y=u_x\)。另一方向的导数满足
	\[
		v_x=-u_y(x,y_0)+\int_{y_0}^{y}u_{xx}(x,t)\,\mathrm{d}t
		=-u_y(x,y_0)-\int_{y_0}^{y}u_{yy}(x,t)\,\mathrm{d}t
		=-u_y(x,y).
	\]
	因此 \(u+\mathrm{i}v\) 满足 Cauchy–Riemann 方程且为 \(C^1\)，从而全纯。这也说明 \(u\) 实际上局部是解析函数的实部，具有任意阶连续偏导。

	若全局的 \(v\) 存在，被积一形式就是 \(\mathrm{d}v\)，其沿闭曲线的积分为零。反过来，假设式\eqref{g4:eq:0103:conjugate-period}成立。固定 \(a\in\Omega\)，对任意 \(z\in\Omega\) 取从 \(a\) 到 \(z\) 的分片 \(C^1\) 道路，定义 \(v(z)\) 为该道路上的积分。开连通的平面区域内任意两点都能用折线连接：从 \(a\) 能以折线到达的点构成非空开闭子集，故就是全域。两条道路的积分之差是拼接闭曲线上的积分，因此 \(v\) 与道路选择无关。把终点在一个小矩形内移动，局部积分公式给出 \(v_x=-u_y\)、\(v_y=u_x\)，所以所得 \(v\) 确为全局调和共轭。两个共轭之差的两个偏导数都为零，沿连接任意两点的折线积分便知其为常数。
\end{proof}

\begin{example}
	\label{g4:exm:0103:log-modulus}
	函数 \(u(z)=\log\lvert z\rvert\) 在 \(\mathbb C^*\) 上调和。写成 \(u(x,y)=\tfrac12\log(x^2+y^2)\)，直接求导得
	\[
		-u_y\,\mathrm{d}x+u_x\,\mathrm{d}y
		=\frac{-y\,\mathrm{d}x+x\,\mathrm{d}y}{x^2+y^2}.
	\]
	在单位圆 \(z=\mathrm{e}^{\mathrm{i}t}\)、\(0\leq t\leq2\pi\) 上，其积分为 \(2\pi\)。因此 \(u\) 没有全局调和共轭。局部共轭是辐角；沿单位圆走一周，辐角增加 \(2\pi\)，正是式\eqref{g4:eq:0103:conjugate-period}检测到的障碍。
\end{example}

调和函数的值还可以从周围的圆周恢复。对 \(u=\operatorname{Re}f\)，这一现象来自 Cauchy 积分公式；以下证明直接从调和方程出发，因而不需要先在整个圆盘上选择共轭。

\begin{proposition}[调和函数的平均值性质]
	\label{g4:prop:0103:harmonic-mean}
	设 \(u\) 是开集 \(U\subset\mathbb C\) 上的实值调和函数，\(\overline{D(a,R)}\subset U\)。则对 \(0<r\leq R\)，有
	\[
		u(a)=\frac1{2\pi}\int_0^{2\pi}u(a+r\mathrm{e}^{\mathrm{i}\theta})\,\mathrm{d}\theta
		=\frac1{\pi r^2}\int_{D(a,r)}u(x,y)\,\mathrm{d}x\,\mathrm{d}y.
	\]
\end{proposition}

\begin{proof}
	令 \(v(r,\theta)=u(a+r\mathrm{e}^{\mathrm{i}\theta})\)，圆周平均为 \(m(r)=(2\pi)^{-1}\int_0^{2\pi}v(r,\theta)\,\mathrm{d}\theta\)。极坐标中的链式法则给出
	\[
		m''(r)+\frac1r m'(r)
		=\frac1{2\pi}\int_0^{2\pi}
		\left((u_{xx}+u_{yy})(a+r\mathrm{e}^{\mathrm{i}\theta})
		-\frac1{r^2}v_{\theta\theta}(r,\theta)\right)
		\,\mathrm{d}\theta=0.
	\]
	这里角向二阶导数的一周积分为零，调和性使 \(u_{xx}+u_{yy}=0\)。因此 \((rm'(r))'=0\)。\(u\) 在中心附近的梯度有界，所以 \(m'(r)\) 有界，\(rm'(r)\rightarrow0\)。于是 \(m'(r)=0\)，再由 \(m(r)\rightarrow u(a)\) 得第一式。用极坐标对圆周平均再乘 \(2r\,\mathrm{d}r\) 积分，得到圆盘平均公式。
\end{proof}

\begin{corollary}[调和函数的极值原理]
	\label{g4:cor:0103:harmonic-maximum}
	设 \(u\) 是区域 \(\Omega\subset\mathbb C\) 上的实值调和函数。若 \(u\) 在内点取得局部最大值或局部最小值，则 \(u\) 为常值。若 \(\Omega\) 有界且 \(u\in C(\overline\Omega)\)，则 \(u\) 的最大值和最小值均在 \(\partial\Omega\) 上取得。特别地，两个在 \(\Omega\) 上调和、在 \(\overline\Omega\) 上连续且边界值相同的实函数必相等。
\end{corollary}

\begin{proof}
	设 \(a\) 为局部最大值点，取足够小的圆盘使其中 \(u\leq u(a)\)。每个以 \(a\) 为中心的小圆周上的平均值都等于 \(u(a)\)，连续非负函数 \(u(a)-u\) 因积分为零而在该圆周恒为零。让半径变化，便知 \(u\) 在一个开圆盘内为常数。

	要把这一结论传到整个连通域，考虑 \(h=u_x-\mathrm{i}u_y\)。调和方程与混合偏导交换分别给出 \(h\) 的两条 Cauchy–Riemann 方程，所以 \(h\) 在 \(\Omega\) 上全纯。它在刚才的圆盘内为零，由\cref{g4:thm:0101:identity}知 \(h\equiv0\)。于是 \(u\) 的梯度在整个区域为零，故 \(u\) 为常数。对 \(-u\) 应用同一论证得到最小值结论。有界域上的边界结论由闭包紧性与内点极值结论得到；最后对两个函数之差应用最大值和最小值结论，便知其恒为零。
\end{proof}

边界条件因此决定了调和延拓的唯一性。存在性需要另行构造：极值原理比较的是已经存在的两个函数。类似地，若全纯函数 \(f\) 在区域内无零点，对 \(1/f\) 应用最大模原理，便得到 \(\lvert f\rvert\) 的最小值原则。无零点这一条件不可省略，例如 \(f(z)=z\) 的模就在原点取得最小值。

\subsection{Schwarz 引理与等号}
\label{g4:subsec:0103:02}
设 \(f\colon\mathbb D\rightarrow\mathbb D\) 全纯并满足 \(f(0)=0\)。由于原点的零点至少有一阶，可以把 \(f(z)\) 除以 \(z\)；对商函数应用最大模原理，就能把值域的界 \(\lvert f\rvert<1\) 改进为随 \(z\) 缩小的界。这是\term[Schwarzyinli]{Schwarz 引理}[Schwarz lemma]的基本想法。\footnote{Hermann Amandus Schwarz（1843--1921），德国数学家，研究共形映射、极小曲面与变分问题。}

\begin{theorem}[Schwarz 引理]
	\label{g4:thm:0103:schwarz}
	若 \(f\colon\mathbb D\rightarrow\mathbb D\) 全纯且 \(f(0)=0\)，则
	\[
		\lvert f(z)\rvert\leq\lvert z\rvert\quad(z\in\mathbb D),
		\qquad \lvert f'(0)\rvert\leq1.
	\]
	若存在非零 \(z_0\in\mathbb D\) 使 \(\lvert f(z_0)\rvert=\lvert z_0\rvert\)，或者 \(\lvert f'(0)\rvert=1\)，则存在 \(\theta\in\mathbb R\) 使 \(f(z)=\mathrm{e}^{\mathrm{i}\theta}z\)。这些旋转也都满足等号。
\end{theorem}

\begin{proof}
	由 \(f\) 在原点的幂级数，函数
	\[
		g(z)=
		\begin{cases}
			f(z)/z,&z\ne0,\\
			f'(0),&z=0
		\end{cases}
	\]
	在整个 \(\mathbb D\) 上全纯。固定 \(0<r<1\)，在 \(\lvert z\rvert=r\) 上有 \(\lvert g(z)\rvert\leq1/r\)。\cref{g4:thm:0103:maximum-modulus}将此界传到闭圆盘 \(\lvert z\rvert\leq r\)。对每个固定 \(z\in\mathbb D\)，令 \(r\uparrow1\)，得到 \(\lvert g(z)\rvert\leq1\)，这同时给出两条不等式。

	任一等号条件都表示 \(\lvert g\rvert\) 在一个内点取得最大值 \(1\)。最大模原理于是使 \(g\) 为模等于 \(1\) 的常数，故 \(g=\mathrm{e}^{\mathrm{i}\theta}\)。反方向直接代入即可。
\end{proof}

等号所给的信息比不等式本身更强：一个非零点的模或原点的导数达到界，就能决定整个函数。若 \(f'(0)=0\)，还可以继续提取零点因子。具体地，若 \(f\) 在原点的零点阶数至少为正整数 \(m\)，对全纯延拓的 \(f(z)/z^m\) 在半径 \(r\) 的圆上估计，再令 \(r\uparrow1\)，得到 \(\lvert f(z)\rvert\leq\lvert z\rvert^m\)。若在非零点取等号，同一极值论证给出 \(f(z)=\mathrm{e}^{\mathrm{i}\theta}z^m\)。例如 \(z^2\) 把靠近原点的点压得更靠近原点，同时保留了一个二阶分歧点。

\ProseParagraphBreak

固定原点的条件使比较特别简洁。若 \(f(0)\ne0\)，\(\lvert f(z)\rvert\leq\lvert z\rvert\) 在 \(z=0\) 便不成立。要对任意全纯圆盘自映射获得同样有效的结论，需要先把任意选定的点及其像同时移到原点。

\subsection{Schwarz–Pick 的度量形式}
\label{g4:subsec:0103:03}
对 \(a\in\mathbb D\)，定义分式线性映射
\begin{equation}
	\label{g4:eq:0103:disk-normalization}
	\phi_a(z)\coloneqq\frac{z-a}{1-\overline a z}.
\end{equation}
分母在 \(\mathbb D\) 内非零，且直接展开分子分母的模平方可得
\begin{equation}
	\label{g4:eq:0103:disk-normalization-identities}
	1-\lvert\phi_a(z)\rvert^2
	=\frac{(1-\lvert a\rvert^2)(1-\lvert z\rvert^2)}{\lvert1-\overline a z\rvert^2},
	\qquad
	\phi_a'(z)=\frac{1-\lvert a\rvert^2}{(1-\overline a z)^2}.
\end{equation}
第一式说明 \(\phi_a\) 把圆盘映入自身；解出 \(z\) 又得到
\[
	\phi_a^{-1}(w)=\frac{w+a}{1+\overline a w}=\phi_{-a}(w).
\]
因此 \(\phi_a\) 是具有全纯逆的圆盘双射，把 \(a\) 移到 \(0\)。

\begin{theorem}[Schwarz--Pick 定理]
	\label{g4:thm:0103:schwarz-pick}
	若 \(f\colon\mathbb D\rightarrow\mathbb D\) 全纯，则对任意 \(z,w\in\mathbb D\)，有
	\begin{equation}
		\label{g4:eq:0103:schwarz-pick-two-points}
		\left\lvert\frac{f(z)-f(w)}{1-\overline{f(w)}f(z)}\right\rvert
		\leq
		\left\lvert\frac{z-w}{1-\overline w z}\right\rvert,
	\end{equation}
	并且对每个 \(z\in\mathbb D\)，有
	\begin{equation}
		\label{g4:eq:0103:schwarz-pick-derivative}
		\frac{\lvert f'(z)\rvert}{1-\lvert f(z)\rvert^2}
		\leq\frac1{1-\lvert z\rvert^2}.
	\end{equation}
	若第一式在某对不同点取等号，或第二式在某一点取等号，则 \(f\) 是圆盘的全纯双射，且逆映射全纯。反过来，每个这样的映射在上述两式中处处取等号。
\end{theorem}

\begin{proof}
	固定 \(a\in\mathbb D\)，把定义域中的 \(a\) 和值域中的 \(f(a)\) 分别移到原点，令
	\[
		g=\phi_{f(a)}\circ f\circ\phi_a^{-1}.
	\]
	这些映射组成交换图
	\[
		\begin{tikzcd}[column sep=large,row sep=large]
			\mathbb D \arrow[r,"f"] \arrow[d,"\phi_a"']
			& \mathbb D \arrow[d,"\phi_{f(a)}"] \\
			\mathbb D \arrow[r,"g"'] & \mathbb D.
		\end{tikzcd}
	\]
	图中 \(a\) 与 \(f(a)\) 均被竖直映射送到 \(0\)，所以 \(g(0)=0\)。由\cref{g4:thm:0103:schwarz}，
	\[
		\lvert\phi_{f(a)}(f(z))\rvert\leq\lvert\phi_a(z)\rvert,
		\qquad \lvert g'(0)\rvert\leq1.
	\]
	第一式正是所需的两点不等式。又由链式法则和式\eqref{g4:eq:0103:disk-normalization-identities}，
	\[
		g'(0)=\frac{1-\lvert a\rvert^2}{1-\lvert f(a)\rvert^2}f'(a),
	\]
	于是得到导数不等式。

	若在不同点 \(z\ne a\) 的两点不等式中取等号，Schwarz 引理的等号条件使 \(g(\zeta)=\mathrm{e}^{\mathrm{i}\theta}\zeta\)；导数等号也给出相同结论。因此
	\[
		f=\phi_{f(a)}^{-1}\circ R_\theta\circ\phi_a,
		\qquad R_\theta(\zeta)=\mathrm{e}^{\mathrm{i}\theta}\zeta.
	\]
	每个因子均有全纯逆，所以 \(f\) 也是这样的双射。反过来，若 \(f\) 及其逆都全纯，分别对二者应用两点不等式，得到互反的两个不等式，故必须取等号。对导数不等式也同样处理，并使用 \((f^{-1})'(f(z))f'(z)=1\)，即得处处等号。
\end{proof}

这个结果称为\term[SchwarzPickdingli]{Schwarz--Pick 定理}[Schwarz--Pick theorem]。\footnote{Georg Alexander Pick（1859--1942），奥地利数学家，研究复变函数、微分方程与几何；格点多边形的面积公式也以其姓氏命名。}例如，当 \(a\ne0\) 时，\(\lvert\phi_a'(a)\rvert=(1-\lvert a\rvert^2)^{-1}>1\)，上述圆盘双射在 \(a\) 附近会放大平面通常的长度。导数不等式中的分母给出了适合圆盘的修正。

\begin{definition}
	\label{g4:def:0103:hyperbolic-distance}
	在单位圆盘上取长度元素
	\[
		\mathrm{d}s_{\mathbb D}=\lambda_{\mathbb D}(z)\lvert\mathrm{d}z\rvert,
		\qquad \lambda_{\mathbb D}(z)\coloneqq\frac2{1-\lvert z\rvert^2}.
	\]
	这称为圆盘的\term[shuangquduliang]{双曲度量}[hyperbolic metric]。分片 \(C^1\) 曲线 \(\gamma\colon[0,1]\rightarrow\mathbb D\) 的长度定义为
	\[
		L_{\mathbb D}(\gamma)=\int_0^1
		\frac{2\lvert\gamma'(t)\rvert}{1-\lvert\gamma(t)\rvert^2}\,\mathrm{d}t.
	\]
	两点间的双曲距离 \(d_{\mathbb D}(z,w)\) 是连接它们的全部这种曲线长度的下确界。
	\symidx{\(\lambda_{\mathbb D}\)：单位圆盘的双曲长度密度}
	\symidx{\(d_{\mathbb D}\)：单位圆盘的双曲距离}
\end{definition}

长度密度处处为正，而且 \(\lambda_{\mathbb D}\geq2\)，所以每条连接 \(z,w\) 的曲线长度至少为 \(2\lvert z-w\rvert\)。圆盘内的直线段给出有限长度的连接。反向参数化保持长度，拼接曲线的长度相加，因而这一构造确实给出有限的距离函数。密度在接近边界时增大，表示靠近单位圆的同样一小段平面长度会被赋予更大的双曲长度。

\begin{proposition}
	\label{g4:prop:0103:hyperbolic-distance-formula}
	对任意 \(z,w\in\mathbb D\)，令
	\[
		q=\left\lvert\frac{z-w}{1-\overline w z}\right\rvert.
	\]
	则 \(0\leq q<1\)，并且
	\begin{equation}
		\label{g4:eq:0103:hyperbolic-distance}
		d_{\mathbb D}(z,w)=\log\frac{1+q}{1-q}
		=2\operatorname{artanh}q.
	\end{equation}
	每个 \(\phi_a\) 和每个旋转都保持这一距离。
\end{proposition}

\begin{proof}
	由式\eqref{g4:eq:0103:disk-normalization-identities}，
	\[
		\lambda_{\mathbb D}(\phi_a(z))\lvert\phi_a'(z)\rvert
		=\lambda_{\mathbb D}(z).
	\]
	链式法则于是说明 \(\phi_a\) 保持每条曲线的双曲长度；它的逆也把连接曲线一一对应起来，所以保持长度的下确界。旋转的长度保持性直接来自模的不变性。取 \(a=w\)，便把待算距离化为原点到某个模为 \(q\) 的点的距离。

	设 \(\gamma\) 从 \(0\) 连到 \(\zeta\)。径向运动给出最短长度的候选；为比较任意曲线，先用光滑近似处理半径在过原点时可能不可微的问题。对 \(\varepsilon>0\)，令
	\[
		r_\varepsilon(t)=\sqrt{\lvert\gamma(t)\rvert^2+\varepsilon^2}-\varepsilon.
	\]
	它分片 \(C^1\)，满足 \(0\leq r_\varepsilon(t)\leq\lvert\gamma(t)\rvert\)，且由求导与复数模的不等式可得
	\[
		\lvert r_\varepsilon'(t)\rvert
		=\frac{\lvert\operatorname{Re}(\overline{\gamma(t)}\gamma'(t))\rvert}
		{\sqrt{\lvert\gamma(t)\rvert^2+\varepsilon^2}}
		\leq\lvert\gamma'(t)\rvert.
	\]
	令
	\[
		A(r)=\log\frac{1+r}{1-r},\qquad A'(r)=\frac2{1-r^2}.
	\]
	曲线像是圆盘内的紧集，所以这里的分母有正下界。由一元微积分的链式法则，
	\[
		L_{\mathbb D}(\gamma)
		\geq\int_0^1\frac{2\lvert r_\varepsilon'(t)\rvert}{1-r_\varepsilon(t)^2}\,\mathrm{d}t
		\geq A(r_\varepsilon(1))-A(r_\varepsilon(0)).
	\]
	因 \(r_\varepsilon(0)=0\)、\(r_\varepsilon(1)\rightarrow\lvert\zeta\rvert\)，令 \(\varepsilon\downarrow0\) 得 \(L_{\mathbb D}(\gamma)\geq A(\lvert\zeta\rvert)\)。径向线段 \(\gamma(t)=t\zeta\) 恰好达到这个下界，因为半径沿曲线单调增加。因此原点到 \(\zeta\) 的距离是 \(A(\lvert\zeta\rvert)\)，从而得到所示公式。
\end{proof}

Schwarz–Pick 的导数不等式现在可以写成
\[
	\lambda_{\mathbb D}(f(z))\lvert f'(z)\rvert
	\leq\lambda_{\mathbb D}(z).
\]
沿任意分片 \(C^1\) 曲线积分，得到 \(L_{\mathbb D}(f\circ\gamma)\leq L_{\mathbb D}(\gamma)\)。在所有连接两点的曲线上取下确界，便有
\begin{equation}
	\label{g4:eq:0103:hyperbolic-contraction}
	d_{\mathbb D}(f(z),f(w))\leq d_{\mathbb D}(z,w).
\end{equation}
因此全纯圆盘自映射都使双曲距离不增加。式\eqref{g4:eq:0103:hyperbolic-distance}中的实函数 \(q\mapsto\log((1+q)/(1-q))\) 严格递增，所以两个不同点间的距离相等，恰好对应\cref{g4:thm:0103:schwarz-pick}的两点等号情形，迫使映射具有全纯逆。

\begin{example}
	\label{g4:exm:0103:square-contraction}
	对 \(f(z)=z^2\)，双曲长度的局部伸缩比为
	\[
		\frac{\lambda_{\mathbb D}(f(z))\lvert f'(z)\rvert}
		{\lambda_{\mathbb D}(z)}
		=\frac{2\lvert z\rvert(1-\lvert z\rvert^2)}{1-\lvert z\rvert^4}
		=\frac{2\lvert z\rvert}{1+\lvert z\rvert^2}<1.
	\]
	最后一步来自 \((1-\lvert z\rvert)^2>0\)。原点处的伸缩比为零，与分歧相符；趋近边界时伸缩比趋于 \(1\)，所以在整个圆盘上不存在统一小于 \(1\) 的伸缩上界。另一方面，原点到实数 \(r\in(0,1)\) 的距离为 \(\log((1+r)/(1-r))\)，当 \(r\uparrow1\) 时趋于无穷。
\end{example}

\subsection{双全纯映射与自同构}
\label{g4:subsec:0103:04}
保持全纯结构的坐标变换要求正向与反向都全纯，这正是\cref{g4:def:0101:biholomorphic}中的双全纯映射。由\cref{g4:thm:0101:open-mapping}及其后的局部逆论证，全纯单射就是到其像的双全纯映射。区域 \(U\subset\mathbb C\) 到自身的双全纯映射称为 \(U\) 的\term[quanchunzitonggou]{全纯自同构}[holomorphic automorphism]，其全体在复合下构成群 \(\operatorname{Aut}(U)\)。对圆盘来说，Schwarz 引理的等号条件足以确定这个群的所有元素。

\begin{theorem}
	\label{g4:thm:0103:disk-aut}
	单位圆盘的全纯自同构恰好是
	\[
		f(z)=\mathrm{e}^{\mathrm{i}\theta}\frac{z-a}{1-\overline a z},
		\qquad a\in\mathbb D,\quad\theta\in\mathbb R.
	\]
	其中 \(a=f^{-1}(0)\) 唯一，\(\theta\) 按模 \(2\pi\) 唯一。特别地，固定原点的全纯自同构都是旋转。
\end{theorem}

\begin{proof}
	所示映射都是式\eqref{g4:eq:0103:disk-normalization}与旋转的复合，已知具有全纯逆。反过来，设 \(f\) 为全纯自同构，取 \(a=f^{-1}(0)\)，令 \(h=f\circ\phi_a^{-1}\)。则 \(h\) 与 \(h^{-1}\) 都固定 \(0\)。对二者分别应用\cref{g4:thm:0103:schwarz}，得到
	\[
		\lvert h(z)\rvert\leq\lvert z\rvert
		=\lvert h^{-1}(h(z))\rvert\leq\lvert h(z)\rvert.
	\]
	于是处处取等号，Schwarz 引理给出 \(h(z)=\mathrm{e}^{\mathrm{i}\theta}z\)，从而 \(f=R_\theta\circ\phi_a\)。\(a\) 由唯一的零点决定；固定 \(a\) 后，旋转因子也由任一非零点的像决定。
\end{proof}

一个圆盘坐标的选择因此有两种自由：先指定哪个点成为原点，再指定该点的方向。固定这两项就能得到唯一性。

\begin{corollary}
	\label{g4:cor:0103:normalized-unique}
	设 \(\Omega\subset\mathbb C\) 是区域，\(a\in\Omega\)。若 \(f,g\colon\Omega\rightarrow\mathbb D\) 均为双全纯映射，且
	\[
		f(a)=g(a)=0,\qquad f'(a)>0,\qquad g'(a)>0,
	\]
	其中导数的不等式表示正实数，则 \(f=g\)。
\end{corollary}

\begin{proof}
	映射 \(h=g\circ f^{-1}\) 是固定 \(0\) 的圆盘自同构，故由\cref{g4:thm:0103:disk-aut}可写成 \(h(z)=\mathrm{e}^{\mathrm{i}\theta}z\)。又
	\[
		\mathrm{e}^{\mathrm{i}\theta}=h'(0)=\frac{g'(a)}{f'(a)}>0.
	\]
	正实数中只有 \(1\) 的模为 \(1\)，所以 \(h\) 是恒等映射，\(f=g\)。
\end{proof}

\begin{example}
	\label{g4:exm:0103:half-plane}
	上半平面 \(\mathbb H=\{w\in\mathbb C:\operatorname{Im}w>0\}\) 与圆盘双全纯等价。具体映射及其逆为
	\[
		\psi(w)=\frac{w-\mathrm{i}}{w+\mathrm{i}},
		\qquad \psi^{-1}(z)=\mathrm{i}\frac{1+z}{1-z}.
	\]
	由
	\[
		1-\lvert\psi(w)\rvert^2
		=\frac{4\operatorname{Im}w}{\lvert w+\mathrm{i}\rvert^2},
		\qquad \lvert\psi'(w)\rvert=\frac2{\lvert w+\mathrm{i}\rvert^2},
	\]
	既可检查 \(\psi(\mathbb H)\subset\mathbb D\)，也可得到圆盘双曲长度的拉回
	\[
		\lambda_{\mathbb D}(\psi(w))\lvert\psi'(w)\rvert\lvert\mathrm{d}w\rvert
		=\frac{\lvert\mathrm{d}w\rvert}{\operatorname{Im}w}.
	\]
	因此上半平面的长度元素可以从圆盘坐标直接恢复。区域在平面中是否有界，不是双全纯等价的不变量。
\end{example}

圆盘自同构的分类解决了圆盘坐标的唯一性问题。对一个给定的平面区域，能否首先找到这样的坐标，需要构造全纯单射并控制其极限；下一节将用正规族把这一存在问题变成一个导数极值问题。
% 纲要条目排为编号小节标题，并进入目录与书签；说明段仍保留于纲要注释中。

% !TeX root = ../../Book4.tex
% 纲要标题：### 1.4 正规族与共形构造
% 纲要位置：几何学计划纲要.md，第 3502 行。
% 正文状态：初稿；尚未正式审核。

% 纲要内容开始
% 
% * 局部有界和 Montel 紧性
% * Hurwitz 定理与单叶极限
% * 平面 Riemann 映照定理的极值证明
% * 边界行为和非单连通域的限制
% 
% ---
% 
% 纲要内容结束

\section{正规族与共形构造}
\label{g4:sec:0104}

把一个平面域变成圆盘，通常找不到可以直接写出的公式。一个可行的办法是先收集所有满足某种归一化条件的单叶映射，再从中选取在指定点放大最多的映射。这立即带来两个问题：趋近最大值的一列映射能否收敛，极限是否还保留单叶性？Cauchy 估计提供前一个问题所需的紧性，零点数的稳定性处理后一个问题。两者合起来便能把一个极值问题变成共形映射的存在证明。

\subsection{局部有界和 Montel 紧性}
\label{g4:subsec:0104:01}

全纯函数在一点附近的变化速度受较大邻域上的函数值控制。因此，对一族全纯函数，统一的局部大小估计还会给出统一的局部连续性估计；这是一般连续函数族没有的性质。

\begin{definition}\label{g4:def:0104:normal-family}
	设 \(\Omega\subset\mathbb C\) 为域。函数列 \(f_j:\Omega\to\mathbb C\) 在每个紧集 \(K\subset\Omega\) 上一致收敛于 \(f\) 时，称其\term[ju bu yi zhi shou lian]{局部一致收敛}[locally uniform convergence]于 \(f\)。对于全纯函数族 \(\mathcal F\subset\mathcal O(\Omega)\)，若每个紧集 \(K\subset\Omega\) 都有常数 \(M_K<\infty\)，使
	\[
		|f(z)|\le M_K\qquad(f\in\mathcal F,\ z\in K),
	\]
	则称 \(\mathcal F\)\term[ju bu you jie]{局部有界}[locally bounded]。若 \(\mathcal F\) 中每个函数列都有一个局部一致收敛于 \(\Omega\) 上全纯函数的子列，则称它在全纯函数空间中为\term[zheng gui zu]{正规族}[normal family]。这里采用有限复值极限；亚纯函数族还可以使用球面距离定义另一种容许无穷远值的正规性。
\end{definition}

局部一致收敛允许靠近边界时失去统一控制。例如 \(f_j(z)=z^j\) 在单位圆盘 \(\mathbb D\) 上局部一致收敛于 \(0\)，因为每个紧集都包含在某个 \(|z|\le r<1\) 中，且 \(|z|^j\le r^j\)。但 \(\sup_{z\in\mathbb D}|z^j|=1\)，所以它们不在整个圆盘上一致收敛。

\begin{proposition}\label{g4:prop:0104:holomorphic-limit}
	设 \(\Omega\subset\mathbb C\) 为域，全纯函数列 \(f_j\) 在 \(\Omega\) 上局部一致收敛于 \(f\)。则 \(f\) 全纯，而且对每个整数 \(k\ge1\)，\(f_j^{(k)}\) 局部一致收敛于 \(f^{(k)}\)。
\end{proposition}
\begin{proof}
	局部一致极限首先是连续函数。取闭圆盘 \(\overline{D(a,R)}\subset\Omega\)。在 \(|z-a|<R\) 上，\cref{g4:thm:0101:cauchy-formula}给出
	\[
		f_j(z)=\frac1{2\pi\mathrm i}\int_{|\zeta-a|=R}
		\frac{f_j(\zeta)}{\zeta-z}\,\mathrm d\zeta.
	\]
	在积分圆周上的一致收敛允许取极限，得到同一公式中的 \(f_j\) 全部换成 \(f\)。对任意 \(0<r<R\)，当 \(|z-a|\le r\) 时，分母的模至少为 \(R-r\)；核函数对 \(z\) 的各阶导数因此都能在积分号下求出。这证明 \(f\) 在圆盘内全纯，并给出
	\[
		\sup_{|z-a|\le r}|f_j^{(k)}(z)-f^{(k)}(z)|
		\le \frac{k!R}{(R-r)^{k+1}}
		\sup_{|\zeta-a|=R}|f_j(\zeta)-f(\zeta)|.
	\]
	右端趋于零。任意紧集可用有限个这样的较小圆盘覆盖，故各阶导数也局部一致收敛。
\end{proof}

要抽取收敛子列，还需要一个紧空间上的基本论证。设 \(K\) 为紧度量空间，连续函数列在 \(K\) 上一致有界且等度连续。取可数稠密集 \(\{x_1,x_2,\ldots\}\subset K\)，依次抽取在 \(x_1,x_2,\ldots\) 处收敛的子列，再取对角子列。它在每个 \(x_\nu\) 处收敛。给定 \(\varepsilon>0\)，等度连续性给出统一的 \(\delta>0\)，使距离小于 \(\delta\) 时函数值相差小于 \(\varepsilon/3\)。用稠密集中的有限多个点组成 \(\delta\)-网；子列在这些点处同时成为 Cauchy 列后，在整个 \(K\) 上的差就小于 \(\varepsilon\)。因此该子列一致收敛。可数稠密集本身可由半径 \(1,1/2,1/3,\ldots\) 的有限网的并得到。

\ProseParagraphBreak

把这个紧性论证与 Cauchy 估计结合，得到\term[Montel ding li]{Montel 定理}[Montel's theorem]。Montel\footnote{Paul Antoine Aristide Montel（1876—1975），法国数学家，发展了全纯函数族与正规族理论。}的结论在这里承担的任务，是把一族函数的大小控制转化为一个极限函数的存在。

\begin{theorem}[Montel]\label{g4:thm:0104:montel}
	设 \(\Omega\subset\mathbb C\) 为域。每个局部有界的全纯函数族 \(\mathcal F\subset\mathcal O(\Omega)\) 都是正规族。
\end{theorem}
\begin{proof}
	先证明任意紧集 \(K\subset\Omega\) 上的等度连续性。取 \(\delta>0\)，使
	\[
		L=\{z\in\mathbb C:\operatorname{dist}(z,K)\le2\delta\}
		\subset\Omega.
	\]
	集合 \(L\) 紧，因而存在对全部 \(f\in\mathcal F\) 通用的上界 \(M_L\)。在距 \(K\) 不超过 \(\delta\) 的点处，半径 \(\delta\) 的圆盘包含在 \(L\) 中，\cref{g4:prop:0101:cauchy-estimate}给出 \(|f'|\le M_L/\delta\)。若 \(z,w\in K\) 且 \(|z-w|<\delta\)，连接二者的线段仍在这个邻域内，故
	\[
		|f(z)-f(w)|
		\le\frac{M_L}{\delta}|z-w|.
	\]
	这就是在 \(K\) 上与 \(f\) 无关的连续性估计。

	现在选择一列紧集穷竭 \(\Omega\)。当 \(\Omega\ne\mathbb C\) 时，可取
	\[
		K_m=\left\{z\in\Omega:|z|\le m,\
		\operatorname{dist}(z,\mathbb C\setminus\Omega)\ge\frac1m\right\},
		\qquad m\ge1;
	\]
	当 \(\Omega=\mathbb C\) 时，取 \(K_m=\{|z|\le m\}\)。这些紧集递增，其内部覆盖 \(\Omega\)，每个 \(\Omega\) 内的紧集都包含在某个 \(K_m\) 中。

	给定 \(\mathcal F\) 中的函数列，利用前述紧空间论证，先抽取在 \(K_1\) 上一致收敛的子列，再从中抽取在 \(K_2\) 上一致收敛的子列，如此继续；空的 \(K_m\) 可略过。取这些嵌套子列的对角子列，它在每个 \(K_m\) 上一致收敛。各个极限在交集上一致，因而定义 \(\Omega\) 上的函数 \(f\)，且所得收敛是局部一致的。由\cref{g4:prop:0104:holomorphic-limit}，\(f\) 全纯。
\end{proof}

局部有界条件必须对函数族统一成立。每个 \(f_j(z)=jz\) 都在紧集上有界，但这列函数在任意含非零点的邻域上没有统一上界，也不可能有局部一致收敛的子列。

\subsection{Hurwitz 定理与单叶极限}
\label{g4:subsec:0104:02}

收敛的全纯函数列仍可能失去某些几何性质。例如 \(f_j(z)=z/j\) 都是单叶函数，极限却是常值。真正需要排除的是这种整体塌缩；一旦极限不恒定，Rouché 定理对零点的控制就会保住单叶性。

\begin{theorem}[Hurwitz]\label{g4:thm:0104:hurwitz}
	设 \(\Omega\subset\mathbb C\) 为域，\(f_j\in\mathcal O(\Omega)\) 局部一致收敛于 \(f\)。若 \(f\) 不恒为零，且 \(a\in\Omega\) 是它的 \(m\) 阶零点，取 \(r>0\)，使 \(\overline{D(a,r)}\subset\Omega\) 且该闭圆盘内没有 \(f\) 的其他零点，则对充分大的 \(j\)，\(f_j\) 在 \(D(a,r)\) 中恰有 \(m\) 个零点，按重数计数。特别地，若每个 \(f_j\) 都无零点，则极限 \(f\) 或者处处非零，或者恒为零。
\end{theorem}
\begin{proof}
	圆周 \(|z-a|=r\) 紧，且其上 \(f\ne0\)，所以
	\[
		\eta=\min_{|z-a|=r}|f(z)|>0.
	\]
	当 \(j\) 足够大时，圆周上有 \(|f_j-f|<\eta\le|f|\)。\cref{g4:thm:0102:rouche}说明 \(f_j\) 与 \(f\) 在圆盘中的零点总数相同，恰为 \(m\)。若所有 \(f_j\) 无零点而非零函数 \(f\) 在一点为零，这个结论就产生矛盾，故得到最后的二择一结论。
\end{proof}

上述结果称为\term[Hurwitz ding li]{Hurwitz 定理}[Hurwitz's theorem]。Hurwitz\footnote{Adolf Hurwitz（1859—1919），德国数学家，研究复函数、代数曲线与数论。}的零点稳定性允许多个零点在极限中合并，但总重数保持不变；例如 \(z^2-j^{-2}\) 的两个简单零点趋向 \(z^2\) 的二重零点。常数列 \(f_j=1/j\) 则说明，处处非零的函数列确实可以趋于恒零函数。

\begin{corollary}\label{g4:cor:0104:univalent-limit}
	设 \(\Omega\subset\mathbb C\) 为域，单叶全纯函数列 \(f_j:\Omega\to\mathbb C\) 局部一致收敛于 \(f\)。则 \(f\) 或为常值函数，或仍为单叶全纯函数。
\end{corollary}
\begin{proof}
	由\cref{g4:prop:0104:holomorphic-limit}，\(f\) 全纯。设 \(f\) 不恒定，固定任意 \(a\in\Omega\)，定义
	\[
		g_j(z)=
		\begin{cases}
			\dfrac{f_j(z)-f_j(a)}{z-a},&z\ne a,\\[6pt]
			f_j'(a),&z=a.
		\end{cases}
	\]
	幂级数说明 \(g_j\) 在 \(a\) 处全纯。由于 \(f_j\) 单叶，\cref{g4:thm:0101:local-power}中局部幂模型的次数只能是 \(1\)，故 \(f_j'(a)\ne0\)；在其他点，单叶性也给出 \(g_j(z)\ne0\)。因此每个 \(g_j\) 在整个 \(\Omega\) 上无零点。

	这些函数局部一致收敛于同样由 \(f\) 定义的差商 \(g\)。在离开 \(a\) 的紧集上，这直接来自 \(f_j\to f\)；在一个以 \(a\) 为中心的较小圆盘上，可用
	\[
		g_j(z)=\int_0^1 f_j'\bigl(a+t(z-a)\bigr)\,\mathrm dt
	\]
	及导数的局部一致收敛验证。由于 \(f\) 不恒定，\(g\) 不恒为零，\cref{g4:thm:0104:hurwitz}于是给出 \(g\) 处处非零。因此 \(z\ne a\) 时 \(f(z)\ne f(a)\)。点 \(a\) 任意，故 \(f\) 单叶。
\end{proof}

在构造映射时，固定一个点处的导数下界就能排除常值极限。这正是下面归一化极值问题的作用。

\subsection{平面 Riemann 映照定理的极值证明}
\label{g4:subsec:0104:03}

圆盘的自同构已经说明：即使一个域与圆盘双全纯，对应也有选择余地。指定一点的像为 \(0\)，再要求该点的导数为正实数，就消除了平移基点和旋转的自由度。以下 \(\mathbb D=\{w\in\mathbb C:|w|<1\}\)，并沿用圆盘自同构
\[
	\phi_b(w)=\frac{w-b}{1-\overline b w},
	\qquad b\in\mathbb D.
\]

\begin{theorem}[平面 Riemann 映照定理]\label{g4:thm:0104:riemann-mapping}
	设 \(\Omega\subsetneq\mathbb C\) 为非空单连通域，\(a\in\Omega\)。则存在唯一的双全纯映射 \(F:\Omega\to\mathbb D\)，满足
	\[
		F(a)=0,\qquad F'(a)>0.
	\]
\end{theorem}

这个结果称为\term[Riemann ying zhao ding li]{Riemann 映照定理}[Riemann mapping theorem]。证明分为三个相互衔接的构造：先找出至少一个从 \(\Omega\) 到圆盘的单叶映射；再在归一化映射中使 \(F'(a)\) 达到上确界；最后证明只要还遗漏圆盘内一点，就能构造导数更大的映射，与极值性矛盾。这一极值方法可与 Ahlfors 的论证对照\cite[第6章第1.1节, 第230--231页]{Ahlfors1979Complex}。

\begin{proof}
	先构造一个有界单叶函数。取 \(b\in\mathbb C\setminus\Omega\)。函数 \(z-b\) 在单连通域 \(\Omega\) 上全纯且处处非零，\cref{g4:prop:0102:logarithm}给出全纯平方根 \(h\)，使
	\[
		h(z)^2=z-b.
	\]
	若 \(h(z_1)=h(z_2)\)，平方后得到 \(z_1=z_2\)，故 \(h\) 单叶。而 \(h(\Omega)\) 与 \(-h(\Omega)\) 不交：若 \(h(z_1)=-h(z_2)\)，平方后仍有 \(z_1=z_2\)，从而 \(h(z_1)=0\)，与 \(z_1\ne b\) 矛盾。

	令 \(c=h(a)\)。由\cref{g4:thm:0101:open-mapping}，存在 \(\rho>0\) 使 \(D(c,\rho)\subset h(\Omega)\)。于是 \(D(-c,\rho)\subset-h(\Omega)\) 完全避开 \(h(\Omega)\)，并且
	\[
		|h(z)+c|\ge\rho\qquad(z\in\Omega).
	\]
	故
	\[
		q(z)=\frac{\rho}{2(h(z)+c)}
	\]
	是模不超过 \(1/2\) 的单叶全纯函数。将 \(q\) 与 \(\phi_{q(a)}\) 复合，再乘一个模为 \(1\) 的常数，就得到一个 \(f:\Omega\to\mathbb D\)，满足 \(f(a)=0\)、\(f'(a)>0\)。这里导数不为零，因为单叶全纯函数的局部幂模型次数为 \(1\)。

	考虑所有满足这些条件的函数组成的集合
	\begin{gather*}
		\mathcal S=\{f:\Omega\to\mathbb D:
		f\text{ 全纯且单叶},\ f(a)=0,\ f'(a)>0\},\\
		M=\sup_{f\in\mathcal S}f'(a).
	\end{gather*}
	刚才的构造说明 \(M>0\)。取 \(r>0\) 使 \(\overline{D(a,r)}\subset\Omega\)，由\cref{g4:prop:0101:cauchy-estimate}有 \(f'(a)\le1/r\)，所以 \(M<\infty\)。选取 \(f_j\in\mathcal S\)，使 \(f_j'(a)\to M\)。函数族一致有界于 \(1\)，\cref{g4:thm:0104:montel}给出局部一致收敛的子列，仍记作 \(f_j\)，其全纯极限记作 \(F\)。

	由\cref{g4:prop:0104:holomorphic-limit}，
	\[
		F(a)=0,\qquad F'(a)=M>0,\qquad |F(z)|\le1.
	\]
	因 \(F\) 不恒定，\cref{g4:thm:0103:maximum-modulus}保证 \(|F(z)|<1\)。再由\cref{g4:cor:0104:univalent-limit}，\(F\) 单叶。因此 \(F\in\mathcal S\)，上确界确实由一个映射达到。

	下面证明 \(F\) 满射。若有 \(\beta\in\mathbb D\setminus F(\Omega)\)，则 \(\beta\ne0\)。将遗漏点移到原点，令
	\[
		H=\phi_\beta\circ F.
	\]
	这是 \(\Omega\) 上处处非零的单叶全纯函数，值域包含于 \(\mathbb D\)。再次使用单连通性和\cref{g4:prop:0102:logarithm}，取全纯函数 \(Q\) 满足 \(Q^2=H\)。于是 \(|Q|<1\)，且 \(Q\) 单叶，因为相等的 \(Q\) 值必给出相等的 \(H\) 值。记 \(\gamma=Q(a)\)，则
	\[
		\gamma^2=-\beta,\qquad 0<|\gamma|<1.
	\]
	映射 \(G=\phi_\gamma\circ Q\) 仍将 \(\Omega\) 单叶地映入 \(\mathbb D\)，且 \(G(a)=0\)。为比较导数，利用
	\[
		|\phi_\beta'(0)|=1-|\beta|^2,
		\qquad |\phi_\gamma'(\gamma)|=\frac1{1-|\gamma|^2},
		\qquad 2\gamma Q'(a)=H'(a),
	\]
	得到
	\begin{align*}
		|G'(a)|
		&=\frac{1-|\beta|^2}{2|\gamma|(1-|\gamma|^2)}F'(a)\\
		&=\frac{1+|\gamma|^2}{2|\gamma|}M
		>M.
	\end{align*}
	最后一个不等式来自 \((1-|\gamma|)^2>0\)。再旋转 \(G\) 的值域，使导数成为正实数，就得到 \(\mathcal S\) 中导数大于 \(M\) 的函数，矛盾。因此 \(F(\Omega)=\mathbb D\)。单叶性使 \(F'\) 处处非零，\cref{g4:thm:0101:inverse}保证其逆映射全纯。

	最后，若 \(F_1,F_2\) 都满足结论，则 \(A=F_2\circ F_1^{-1}\) 是固定 \(0\) 的圆盘自同构。由\cref{g4:thm:0103:disk-aut}，\(A(w)=e^{\mathrm i\theta}w\)。而
	\[
		A'(0)=\frac{F_2'(a)}{F_1'(a)}>0,
	\]
	故 \(e^{\mathrm i\theta}=1\)，从而 \(F_1=F_2\)。
\end{proof}

平方根在证明中出现两次，解决的是两个具体问题。第一次利用域外点，把原域映到一个避开整块圆盘的区域，随后取倒数得到有界函数。第二次利用像域内的遗漏点，在保持单叶性的同时扩大基点处的导数。单连通性保证这两次平方根都能在整个 \(\Omega\) 上选成单值全纯函数；正规族紧性则保证最大化过程不会失去极限。

\subsection{边界行为和非单连通域的限制}
\label{g4:subsec:0104:04}

Riemann 映照定理给出开域之间的双全纯对应。若一个点列在源域中不断逼近边界，它在圆盘中的像也必须离开每个内部紧集。这个结论只需逆映射的连续性，已经能够说明许多边界现象。

\begin{proposition}\label{g4:prop:0104:boundary-escape}
	设 \(\Omega\subset\mathbb C\) 为域，\(F:\Omega\to\mathbb D\) 为双全纯映射，\(z_j\in\Omega\)。则以下两件事等价：点列 \(z_j\) 最终离开 \(\Omega\) 的每个紧集；\(|F(z_j)|\to1\)。特别地，若 \(z_j\to p\in\partial\Omega\)，或 \(|z_j|\to\infty\)，则 \(|F(z_j)|\to1\)。
\end{proposition}
\begin{proof}
	若 \(|F(z_j)|\) 不趋于 \(1\)，则存在 \(r<1\) 和一个子列，其像都位于 \(\overline{D(0,r)}\)。逆映射连续，故
	\[
		F^{-1}\bigl(\overline{D(0,r)}\bigr)
	\]
	是 \(\Omega\) 中的紧集，与点列最终离开每个紧集矛盾。反过来，对任意紧集 \(K\subset\Omega\)，\(F(K)\) 紧且包含在 \(\mathbb D\) 中，故 \(\max_K|F|<1\)。若 \(|F(z_j)|\to1\)，则充分大的 \(j\) 都有 \(z_j\notin K\)。最后两种点列都最终离开 \(\Omega\) 的每个紧集，因而满足结论。
\end{proof}

上式控制的是像点到圆周的径向距离。若要在每个边界点赋予唯一的圆周像，还需要研究从不同方向逼近同一点时的行为。狭缝域给出一个具体区别。

\begin{example}\label{g4:exm:0104:slit-boundary}
	考虑
	\[
		\Omega=\mathbb D\setminus[0,1).
	\]
	每个点有唯一表示 \(z=re^{\mathrm i\theta}\)，其中 \(0<r<1\)、\(0<\theta<2\pi\)。映射 \(z\mapsto(\log r,\theta)\) 把 \(\Omega\) 同胚到凸集 \((-\infty,0)\times(0,2\pi)\)，所以 \(\Omega\) 单连通；\cref{g4:thm:0104:riemann-mapping}给出它到圆盘的双全纯映射。

	然而，这个映射不可能延拓为 Euclid 闭包之间的同胚。因为
	\[
		\overline\Omega=\overline{\mathbb D},
		\qquad
		\overline\Omega\setminus\Omega=\partial\mathbb D\cup[0,1).
	\]
	若存在延拓同胚，它须把 \(\partial\mathbb D\cup[0,1)\) 同胚地映到 \(\partial\mathbb D\)。从源边界删去狭缝上的一点 \(p\in(0,1)\) 后，该集合分成包含 \(0\) 的一段和其余部分；圆周删去一点却仍然连通，故这样的边界同胚不存在。

	局部公式 \(\sqrt z=\sqrt r\,e^{\mathrm i\theta/2}\) 更直观地显示了狭缝的两侧：从上侧逼近 \(p>0\) 时，它趋向 \(\sqrt p\)；从下侧逼近时，它趋向 \(-\sqrt p\)。源域边界中的同一个 Euclid 点携带着两种不同的内部逼近方向。处理边界对应时，必须记录这种区别。
\end{example}

单连通假设也有可计算的作用。设 \(0<r<1\)，圆环
\[
	A_r=\{z\in\mathbb C:r<|z|<1\}
\]
上的坐标函数 \(z\) 没有全纯对数：若 \(e^{L(z)}=z\)，求导得 \(L'(z)=1/z\)，但沿 \(|z|=\rho\)、\(r<\rho<1\) 积分会同时给出
\[
	\int_{|z|=\rho}L'(z)\,\mathrm dz=0,
	\qquad
	\int_{|z|=\rho}\frac{\mathrm dz}{z}=2\pi\mathrm i.
\]
若有双全纯映射 \(F:A_r\to\mathbb D\)，其逆映射 \(G:\mathbb D\to A_r\) 是处处非零的全纯函数。由\cref{g4:prop:0102:logarithm}，\(G\) 有全纯对数 \(H\)，于是 \(H\circ F\) 就是 \(A_r\) 上坐标函数的全纯对数，产生矛盾。因此圆环不能双全纯地变成圆盘。

\ProseParagraphBreak

另一个被排除的域是整个 \(\mathbb C\)。若 \(F:\mathbb C\to\mathbb D\) 全纯，则在任意 \(a\in\mathbb C\) 处，半径 \(R\) 的 Cauchy 估计给出 \(|F'(a)|\le1/R\)。令 \(R\to\infty\) 得 \(F'(a)=0\)，所以 \(F\) 恒定。这说明定理中的“单连通真子域”同时包含两个作用不同的条件：单连通性消除选取全纯根时的环路障碍，域外点则启动有界单叶函数的构造。


\begin{chaptersummary}
	全纯性首先是一个一阶条件。三角形积分的消去和 Cauchy 公式使这个条件产生全部高阶导数与幂级数，从而把非恒定全纯函数在每点附近化为 \(w\mapsto w^m\)。整数 \(m\) 既是零点重数，也是邻近原像的数目；留数和辐角原理把这些局部整数合成边界上的绕数。Rouché 定理说明，只要边界不经过零，这个计数便在适当扰动中保持。

	\ProseParagraphBreak

	解析延拓的对象是带中心的函数芽。局部相等保证接续唯一，沿所有道路可延拓与同伦不变性进一步给出单连通域上的单值函数。极值原理把值域的控制转化为刚性，Schwarz--Pick 将其写成双曲长度与距离不增加的形式。正规族则利用同一类积分估计，把局部有界性转成紧性；零点稳定性使非退化的单叶极限继续单叶。对归一化映射的导数取极值，再以遗漏值构造更优候选，就得到平面 Riemann 映照定理。所得坐标首先定义在域内，边界的逐点对应另受边界几何的约束。
\end{chaptersummary}

\BookExercises

\begin{exercise}
	\label{g4:ex:0101:rigid-images}
	设 \(\Omega\subset\mathbb C\) 是区域，\(f\in\mathcal O(\Omega)\)。分别证明：若 \(f\) 的值都落在同一条实仿射直线上，或者 \(|f|\) 恒定，则 \(f\) 为常数。第一种情形另用 Cauchy--Riemann 方程证明。
\end{exercise}
\begin{solution}
	直线和圆都不包含平面开集，所以\cref{g4:thm:0101:open-mapping}排除了非常值函数。若常模为零，结论由定义直接成立。

	对直线情形，先用非零复数的乘法和一次平移把该直线送到实轴，得到全纯实值函数 \(g=u+\mathrm i0\)。Cauchy--Riemann 方程给出 \(u_x=u_y=0\)。区域内任意两点能以有限条直线段相连，沿这些线段积分得到 \(u\) 为常数，还原仿射变换即可。
\end{solution}

\begin{exercise}
	\label{g4:ex:0101:polynomial-growth}
	设整函数 \(f\) 满足 \(|f(z)|\le C(1+|z|)^m\)，其中 \(C>0\)、\(m\) 为非负整数。证明 \(f\) 是次数至多为 \(m\) 的多项式。再用有界整函数必为常数这一结论，证明每个非常数复多项式至少有一个复根。
\end{exercise}
\begin{solution}
	对 \(n>m\)，以原点为中心用\cref{g4:prop:0101:cauchy-estimate}，得到
	\[
		|f^{(n)}(0)|\le n!C\frac{(1+R)^m}{R^n}\longrightarrow0\qquad(R\to\infty).
	\]
	因此所有 \(n>m\) 的幂级数系数为零。任意有限圆盘上都可用\cref{g4:thm:0101:cauchy-formula}展开，同一组系数便在整个平面表示 \(f\)，所以它是所述多项式。

	若非常数多项式 \(P\) 处处非零，则 \(1/P\) 整。最高次项表明 \(|P(z)|\to\infty\) 当 \(|z|\to\infty\)，故 \(1/P\) 在某大圆盘外有界；在闭圆盘内又因连续且分母非零而有界。因此 \(1/P\) 为常数，矛盾。
\end{solution}

\begin{exercise}
	\label{g4:ex:0101:composition-order}
	设 \(f\) 在 \(a\) 附近全纯，\(g\) 在 \(b=f(a)\) 附近全纯，二者均不局部恒定。记 \(m=\operatorname{ord}_a(f-b)\)、\(n=\operatorname{ord}_b(g-g(b))\)。证明复合映射在 \(a\) 的局部重数为 \(mn\)。并证明重数不随源、靶的双全纯局部坐标改变。
\end{exercise}
\begin{solution}
	写 \(f(z)-b=(z-a)^m u(z)\)、\(g(w)-g(b)=(w-b)^n v(w)\)，其中 \(u(a)v(b)\ne0\)。代入得
	\[
		g(f(z))-g(b)=(z-a)^{mn}u(z)^n v(f(z)).
	\]
	最后一个因子在 \(a\) 非零，故阶数恰为 \(mn\)。双全纯局部坐标及其逆的导数非零，局部重数都是 \(1\)；前后复合只将重数乘以 \(1\)，于是得到坐标无关性。
\end{solution}

\begin{exercise}
	\label{g4:ex:0101:local-root-radius}
	在原点附近取满足 \(s(0)=1\)、\(s(z)^2=1+z\) 的全纯分支。求其幂级数的前三个非恒定项，并证明以原点为中心的收敛半径恰为 \(1\)。
\end{exercise}
\begin{solution}
	平方映射在 \(1\) 的导数为 \(2\)，逆函数定理保证该分支存在且由初值唯一。将 \(s=1+a_1z+a_2z^2+a_3z^3+\cdots\) 平方比较系数，得到 \(2a_1=1\)、\(2a_2+a_1^2=0\)、\(2a_3+2a_1a_2=0\)，所以
	\[
		s(z)=1+\frac z2-\frac{z^2}8+\frac{z^3}{16}+\cdots.
	\]
	在 \(\mathbb D\) 上，函数 \(1+z\) 不为零，由\cref{g4:prop:0102:logarithm}可取与该初值一致的全纯平方根。故幂级数在每个半径小于 \(1\) 的圆盘收敛，收敛半径至少为 \(1\)。若大于 \(1\)，级数和就在 \(-1\) 附近全纯。恒等定理使 \(s^2=1+z\) 仍成立，于是 \(s(-1)=0\)；对等式求导又要求 \(2s(-1)s'(-1)=1\)，矛盾。
\end{solution}

\begin{exercise}
	\label{g4:ex:0101:residue-order}
	设 \(a\in\mathbb C\setminus\{0\}\)，\(F(z)=(z^2-a^2)/z^3\)。计算 \(F'/F\) 的全部留数，并求当 \(r>0\) 且 \(r\ne|a|\) 时，闭曲线 \(F(r e^{\mathrm i t})\)、\(0\le t\le2\pi\)，绕原点的绕数。解释为何本例可以出现负绕数。
\end{exercise}
\begin{solution}
	对数导数为
	\[
		\frac{F'}F=\frac1{z-a}+\frac1{z+a}-\frac3z.
	\]
	在 \(a,-a,0\) 的留数依次为 \(1,1,-3\)，没有其他极点。由\cref{g4:thm:0102:argument}，\(0<r<|a|\) 时绕数为 \(-3\)，\(r>|a|\) 时为 \(1+1-3=-1\)。当 \(r=|a|\) 时像曲线经过零，绕数不定义。负号来自圆盘内部的三阶极点；对于在整个闭圆盘邻域全纯的函数，辐角原理只有零点贡献，不会产生这种负计数。
\end{solution}

\begin{exercise}
	\label{g4:ex:0101:rouche-polynomial}
	不求出各根的显式表达式，证明 \(P(z)=z^6+4z^2+1\) 在 \(|z|<1\) 内恰有两个根，在 \(1<|z|<2\) 内恰有四个根，均按重数计算。并证明全部根都是单根。
\end{exercise}
\begin{solution}
	在 \(|z|=1\) 上，\(|z^6+1|\le2<4=|4z^2|\)。Rouché 定理给出与 \(4z^2\) 相同的两个零点，且边界无零点。在 \(|z|=2\) 上，\(|4z^2+1|\le17<64=|z^6|\)，所以大圆盘内有六个零点、边界仍无零点。相减得到环域中的四个。

	若 \(z\) 是重根，则 \(P'(z)=2z(3z^4+4)=0\)。因 \(P(0)=1\)，必有 \(z\ne0\)、\(z^4=-4/3\)。代回 \(P(z)=0\) 得 \(8z^2/3+1=0\)，即 \(z^2=-3/8\)。于是 \(z^4=9/64\)，与 \(z^4=-4/3\) 矛盾，所以不存在重根。
\end{solution}

\begin{exercise}
	\label{g4:ex:0101:annular-roots}
	设 \(A=\{r<|z|<R\}\)，其中 \(0<r<R<\infty\)，\(k\in\mathbb Z\)，\(n\ge2\)。判断 \(z^k\) 在 \(A\) 上何时有全纯对数，何时有全纯 \(n\) 次根，并给出充分性构造。
\end{exercise}
\begin{solution}
	取 \(r<\rho<R\)，沿 \(|z|=\rho\) 逆时针积分。若 \(e^L=z^k\)，则 \(L'=k/z\)，所以
	\[
		0=\int_{|z|=\rho}L'(z)\,\mathrm{d}z=2\pi\mathrm i k.
	\]
	必要条件为 \(k=0\)，此时取 \(L=0\) 即可。

	若 \(g^n=z^k\)，则 \(g\) 处处非零，微分得到 \(n g'/g=k/z\)。积分并除以 \(2\pi\mathrm i\)，左端是整数绕数的 \(n\) 倍，所以必须有 \(n\mid k\)。反过来，若 \(k=n\ell\)，取 \(g(z)=z^\ell\) 即可，负的 \(\ell\) 在环域上同样全纯。任意两个根之比取值于有限个 \(n\) 次单位根，连通性使其恒定，故全部根正是这些单位根乘以 \(z^\ell\)。
\end{solution}

\begin{exercise}
	\label{g4:ex:0101:cubic-monodromy}
	从 \(z=1\) 处取值为 \(1\) 的 \(z^{1/3}\) 函数芽出发，沿 \(\gamma(t)=e^{2\pi\mathrm i t}\)、\(0\le t\le1\)，作解析延拓。求终点函数芽，并求绕行三周后的结果。说明为何在局部处处存在逆分支，仍不能在 \(\mathbb C^*\) 上得到单值立方根。
\end{exercise}
\begin{solution}
	沿路径取值 \(w(t)=e^{2\pi\mathrm i t/3}\)。它满足 \(w(t)^3=\gamma(t)\)，且 \(w(t)\ne0\)。立方映射在各 \(w(t)\) 处导数 \(3w(t)^2\) 非零，给出相应的全纯局部逆；紧参数区间可取有限个这样的邻域组成延拓链。终点值为 \(e^{2\pi\mathrm i/3}\)，所以终点芽等于起点芽乘以该常数；此结论也由局部根的唯一性得到。三周之后乘子为 \(e^{2\pi\mathrm i}=1\)，回到初始芽。一次环路的返回已经改变芽，故不能来自 \(\mathbb C^*\) 上的单值全纯函数；上一题的 \(k=1,n=3\) 给出同一障碍的积分表达。
\end{solution}

\begin{exercise}
	\label{g4:ex:0101:two-fixed-points}
	设 \(f:\mathbb D\to\mathbb D\) 全纯。证明：若 \(f\) 固定两个不同的内点，则 \(f\) 是恒等映射；若 \(f(a)=a\) 且 \(f'(a)=1\)，也有同一结论。并给出只固定一个内点而不恒等的例子。
\end{exercise}
\begin{solution}
	令 \(h=\phi_a\circ f\circ\phi_a^{-1}\)，其中 \(\phi_a(z)=(z-a)/(1-\overline a z)\)。它固定 \(0\)。若另有不同固定点 \(b\)，则 \(h\) 还固定非零点 \(\phi_a(b)\)。Schwarz 引理的等号情形使 \(h(z)=e^{\mathrm i\theta}z\)，固定该非零点又使 \(e^{\mathrm i\theta}=1\)。若改用导数条件，则链式法则与 \(f(a)=a\) 给出 \(h'(0)=f'(a)=1\)，同样由等号条件知 \(h\) 为恒等。还原共轭得到 \(f\) 为恒等。旋转 \(f(z)=-z\) 固定原点但并非恒等，说明一个固定点本身不够。
\end{solution}

\begin{exercise}
	\label{g4:ex:0101:harmonic-oscillation}
	设 \(0<r<R\)，\(c_0,c_1\in\mathbb R\)。在圆环 \(A=\{r<|z|<R\}\) 上构造一个调和函数，使它连续延拓到闭圆环，在内圆周取常值 \(c_0\)，在外圆周取常值 \(c_1\)。证明解唯一，并判断何时存在全局调和共轭。
\end{exercise}
\begin{solution}
	调和函数 \(\log|z|\) 在两条圆周上分别取 \(\log r\)、\(\log R\)。因此可作仿射组合
	\[
		u(z)=c_0+\frac{c_1-c_0}{\log(R/r)}\log\frac{|z|}{r}.
	\]
	它调和、连续至边界，并满足所给边值。若有另一个解，两解之差在边界为零，\cref{g4:cor:0103:harmonic-maximum}使其在整个圆环为零，故解唯一。

	令 \(A_0=(c_1-c_0)/\log(R/r)\)。局部调和共轭为 \(A_0\arg z\) 加常数，其微分沿任意正向同心圆的积分为 \(2\pi A_0\)。由\cref{g4:prop:0103:harmonic-conjugate}，全局共轭存在必有 \(A_0=0\)，即 \(c_1=c_0\)。反之，此时 \(u\) 为常数，取任意实常数作其共轭即可。
\end{solution}

\begin{exercise}
	\label{g4:ex:0101:degenerate-limits}
	设 \(\Omega\subset\mathbb C\) 是区域，单叶全纯函数列 \(f_n\) 局部一致收敛于非恒定函数 \(f\)。证明：每个紧集 \(K\subset f(\Omega)\) 都从某一项起包含在 \(f_n(\Omega)\) 内，而且逆映射 \(f_n^{-1}\) 在 \(K\) 上一致趋于 \(f^{-1}\)。
\end{exercise}
\begin{solution}
	由\cref{g4:cor:0104:univalent-limit}，\(f\) 单叶，所以 \(g=f^{-1}\) 有定义。固定 \(w_0\in f(\Omega)\) 及 \(\varepsilon>0\)，取 \(0<\rho<\varepsilon/2\)，使 \(\overline{D(g(w_0),\rho)}\subset\Omega\)。在该圆周上 \(f-w_0\) 无零点，令其模的最小值为 \(\eta>0\)。

	对 \(|w-w_0|<\eta/3\)，Rouché 定理说明 \(f-w\) 在此圆盘中仍恰有一个零点，因此 \(w\in f(\Omega)\)，且 \(g(w)\) 落在圆盘内。又因圆周上 \(f_n\to f\) 一致，充分大的 \(n\) 满足 \(|f_n-f|<\eta/3\)。这时
	\[
		|(f_n-w)-(f-w_0)|<2\eta/3<|f-w_0|
	\]
	在圆周上成立，故 \(f_n-w\) 在圆盘内恰有一个零点。单叶性说明它就是 \(f_n^{-1}(w)\)，于是
	\[
		|f_n^{-1}(w)-g(w)|<2\rho<\varepsilon
	\]
	在整个 \(D(w_0,\eta/3)\) 上成立。

	对 \(K\) 中各点作此构造，取有限个目标圆盘覆盖 \(K\)，再取有限个指标阈值的最大值。由此同时得到 \(K\subset f_n(\Omega)\) 以及 \(K\) 上误差小于 \(\varepsilon\)，证明了所需一致收敛。
\end{solution}

\begin{exercise}
	\label{g4:ex:0101:conformal-radius}
	设 \(\Omega\subsetneq\mathbb C\) 为单连通区域，\(a\in\Omega\)，\(F:\Omega\to\mathbb D\) 为满足 \(F(a)=0\)、\(F'(a)>0\) 的双全纯映射。定义 \(r_\Omega(a)=1/F'(a)\)。证明：若 \(\Omega_1\subset\Omega_2\subsetneq\mathbb C\) 都是单连通区域且 \(a\in\Omega_1\)，则 \(r_{\Omega_1}(a)\le r_{\Omega_2}(a)\)，等号当且仅当两域相同。计算 \(\Omega=D(c,R)\) 时的数值。
\end{exercise}
\begin{solution}
	Riemann 映照定理与归一化唯一性保证定义存在且不依赖选择。取两域的归一化映射 \(F_1,F_2\)。函数 \(H=F_2\circ F_1^{-1}:\mathbb D\to\mathbb D\) 全纯且固定 \(0\)，其中使用包含映射将 \(\Omega_1\) 看作 \(\Omega_2\) 的子域。于是
	\[
		0<H'(0)=\frac{F_2'(a)}{F_1'(a)}
		=\frac{r_{\Omega_1}(a)}{r_{\Omega_2}(a)}\le1.
	\]
	等号时 Schwarz 引理给出 \(H\) 为恒等，因此 \(F_2(\Omega_1)=\mathbb D\)。\(F_2\) 在 \(\Omega_2\) 上单射且已满射到圆盘，故 \(\Omega_1=\Omega_2\)。反向由定义立即成立。

	令 \(\alpha=(a-c)/R\)，则 \(F(z)=\phi_\alpha((z-c)/R)\) 满足所需归一化，且
	\[
		F'(a)=\frac1{R(1-|\alpha|^2)},\qquad
		r_{D(c,R)}(a)=\frac{R^2-|a-c|^2}{R}.
	\]
	这也显示该量依赖基点：在圆心为 \(R\)，靠近圆周时趋于零。
\end{solution}

\begin{exercise}
	\label{g4:ex:0101:transport-metric}
	设 \(\Omega\subsetneq\mathbb C\) 为单连通区域，\(F:\Omega\to\mathbb D\) 为任意双全纯映射。证明
	\[
		\lambda_\Omega(z)=\frac{2|F'(z)|}{1-|F(z)|^2}
	\]
	与 \(F\) 的选择无关。若 \(G:\Omega_1\to\Omega_2\) 是两个这样的区域之间的全纯映射，证明 \(\lambda_{\Omega_2}(G(z))|G'(z)|\le\lambda_{\Omega_1}(z)\)，并说明某点等号的含义。
\end{exercise}
\begin{solution}
	另一双全纯坐标 \(\widetilde F\) 与 \(F\) 相差圆盘自同构 \(A=\widetilde F\circ F^{-1}\)。\cref{g4:thm:0103:schwarz-pick}的等号给出
	\[
		\frac{|A'(w)|}{1-|A(w)|^2}=\frac1{1-|w|^2}.
	\]
	令 \(w=F(z)\) 并乘以 \(2|F'(z)|\)，即得选择无关性。

	取两域的圆盘坐标 \(F_1,F_2\)，对 \(F_2\circ G\circ F_1^{-1}\) 用 Schwarz--Pick 导数不等式，再按链式法则展开，得到所求式。若某点取等号，该圆盘自映射便是自同构，所以 \(G\) 是两域之间的双全纯映射。反之双全纯映射使该不等式处处为等号。
\end{solution}

\begin{exercise}
	\label{g4:ex:0101:boundary-escape}
	设 \(\Omega\subsetneq\mathbb C\) 是单连通区域，以前一题的密度 \(\lambda_\Omega\) 定义曲线长度及距离 \(d_\Omega\)。证明任意双全纯映射 \(F:\Omega\to\mathbb D\) 都是这两个长度距离之间的等距映射，并证明 \((\Omega,d_\Omega)\) 完备。对固定 \(a\in\Omega\)，证明 \(d_\Omega(a,z_j)\to\infty\) 当且仅当点列 \(z_j\) 最终离开 \(\Omega\) 的每个紧集。
\end{exercise}
\begin{solution}
	由密度的定义及链式法则，每条分片 \(C^1\) 曲线满足 \(L_\Omega(\gamma)=L_{\mathbb D}(F\circ\gamma)\)。\(F\) 与其逆把连接两点的曲线一一对应，取下确界得
	\[
		d_\Omega(z,w)=d_{\mathbb D}(F(z),F(w)).
	\]
	所以只需在圆盘上证明完备性。设 \((w_j)\) 是双曲距离的 Cauchy 列。由 \(d_{\mathbb D}(z,w)\ge2|z-w|\)，它在平面中也是 Cauchy 列，趋于某个 \(w\in\overline{\mathbb D}\)。另一方面，Cauchy 条件及三角不等式使 \(d_{\mathbb D}(0,w_j)\) 有界。公式\eqref{g4:eq:0103:hyperbolic-distance}给出
	\[
		d_{\mathbb D}(0,w_j)=\log\frac{1+|w_j|}{1-|w_j|},
	\]
	所以 \(|w_j|\) 有一个严格小于 \(1\) 的统一上界，极限 \(w\) 在圆盘内。在 \(w\) 的一个小闭圆盘内，长度密度有界；用连接 \(w_j\) 与 \(w\) 的直线段估计长度，便有 \(d_{\mathbb D}(w_j,w)\to0\)。这证明完备性。

	最后选择满足 \(F(a)=0\) 的圆盘坐标。上式说明 \(d_\Omega(a,z_j)\to\infty\) 等价于 \(|F(z_j)|\to1\)，后者由\cref{g4:prop:0104:boundary-escape}等价于离开源域的每个紧集。通常平面距离有限的边界，在这个长度结构下也位于无穷远处。
\end{solution}
